\documentclass[11pt]{amsart}
\usepackage[margin=1in]{geometry}
\usepackage{amsmath,amssymb,amsthm,mathtools}
\usepackage{booktabs}
\usepackage{enumitem}
\usepackage{xcolor}
\usepackage{hyperref}
\usepackage{microtype}
\usepackage[T1]{fontenc}
\usepackage{lmodern}

\newtheorem{theorem}{Theorem}[section]
\newtheorem{proposition}[theorem]{Proposition}
\newtheorem{lemma}[theorem]{Lemma}
\newtheorem{corollary}[theorem]{Corollary}
\theoremstyle{definition}
\newtheorem{definition}[theorem]{Definition}
\newtheorem{remark}[theorem]{Remark}

\newcommand{\eps}{\varepsilon}
\newcommand{\Fcal}{\mathcal F}
\newcommand{\Pcal}{\mathcal P}
\newcommand{\Dcal}{\mathcal D}
\newcommand{\E}{\mathbb E}
\newcommand{\Prb}{\mathbb P}
\newcommand{\R}{\mathbb R}

\hypersetup{
 colorlinks=true,
 linkcolor=blue!50!black,
 citecolor=blue!50!black,
 urlcolor=blue!50!black,
 pdftitle={The Annealed Critical Window for Growing-Radius Domination in Random Regular Graphs},
 pdfauthor={Levi Neuwirth},
 pdfsubject={Growing-radius domination and random regular graphs},
 pdfkeywords={random regular graphs, domination, configuration model, method of types, first moment}
}

\title[The Annealed Critical Window for Growing-Radius Domination]{The Annealed Critical Window for\\Growing-Radius Domination in Random Regular Graphs}
\author{Levi Neuwirth}
\address{Brown University}
\date{July 24, 2026}
\keywords{random regular graphs, domination, configuration model, method of types, first moment}

\begin{document}

\begin{abstract}
Fix $d\ge3$ and let
\[
 B_h=1+d\frac{(d-1)^h-1}{d-2}
\]
be the radius-$h$ volume of the infinite $d$-regular tree.  We determine the bounded annealed critical window for distance-$h$ domination in the random $d$-regular configuration model.  Writing $L_h=\log B_h$, uniformly for bounded $s$ we prove
\[
 \frac{B_h}{L_h^2}
 \Psi_{d,h}\!\left(
 \frac{L_h-2\log L_h+s}{B_h}
 \right)
 \longrightarrow 1-e^{-s},
\]
where $\Psi_{d,h}$ is the exact microcanonical first-moment exponent.  Equivalently,
\[
 H(\alpha)-\Psi_{d,h}(\alpha)
 =(1-\alpha)^{B_h}(1+o(1))
\]
throughout the critical window.  The lower zero is pinned within
$B_h^{-1/7+o(1)}$ of the scalar coupon root
$H(C/B_h)=(1-C/B_h)^{B_h}$ and therefore has a complete fixed-order inverse-logarithmic expansion; in particular,
\[
 B_h\alpha_h^{\rm ann}
 =L_h-2\log L_h
 +\frac{3\log L_h-1}{L_h}
 +O\!\left(\frac{(\log L_h)^2}{L_h^2}\right).
\]
Consequently, for every fixed $\omega>0$ and the stated growth condition, a uniformly random simple $d$-regular graph satisfies
\[
 \gamma_h(G_{n,d})\ge
 \frac n{B_h}
 \bigl(L_h-2\log L_h-\omega\bigr)
\]
with high probability.

The proof uses exact graph-distance labels, a two-sided method-of-types count, and an exact reduction to a compact tridiagonal variational functional.  Although that functional is nonconcave, boundary repulsion and a monotone reverse transfer identify its unique global optimizer.  Quantitative stable/free shadowing gives the activity law at the required scale.  The new upper half of the critical-window argument comes from an explicit capped-free profile: all nonterminal layers retain the free Bernoulli entropy, and the entire entropy defect is the cost of one terminal nonemptiness conditioning event.  The result also transfers as a lower bound to internally two-path $(h,2)$ domination.  It remains annealed: quenched matching and the direct two-branch leading constant are open.
\end{abstract}

\maketitle

\noindent\textbf{Keywords.} random regular graphs; domination number; growing-radius domination; configuration model; method of types; belief propagation; first moment.

\section{Introduction and related work}\label{sec:intro}

\subsection{Main result and scale}
A set $S\subseteq V(G)$ is \emph{distance-$h$ dominating} if every vertex of $G$ lies within graph distance $h$ of $S$; its minimum size is denoted $\gamma_h(G)$.  For graphs of maximum degree $d$, one selected vertex can cover at most
\[
 B_h=1+d\frac{(d-1)^h-1}{d-2}
\]
vertices, so the elementary volume bound is $\gamma_h(G)\ge n/B_h$.  In ideal tree-ball geometry, independent selection and patching place the natural covering scale at $n\log B_h/B_h$.

The first-moment balance is subtler than the leading scale.  If $\alpha=C/B_h$, subset entropy is approximately
$C(\log B_h-\log C+1)/B_h$, while the probability that a tree ball is missed is approximately $e^{-C}$.  Equating these terms predicts
\[
 C=\log B_h-2\log\log B_h+O(1).
\]
The present paper determines that bounded annealed window and its entire limiting profile.  With $L_h=\log B_h$, uniformly for bounded $s$,
\[
 \frac{B_h}{L_h^2}
 \Psi_{d,h}\!\left(
 \frac{L_h-2\log L_h+s}{B_h}
 \right)
 \longrightarrow 1-e^{-s}.
\]
Thus the annealed exponent changes sign at $s=0$, and its lower zero is governed, to power accuracy in $B_h$, by the scalar equation
\[
 H(C/B_h)=(1-C/B_h)^{B_h}.
\]
This gives not only the first two terms but the complete fixed-order inverse-logarithmic expansion of the annealed transition.

For the random graph itself, the first moment yields a fixed-slack lower theorem: for every fixed $\omega>0$, subject to
\[
 \frac{n(\log B_h)^2}{B_h}\gg h\log n,
\]
one has with high probability
\[
 \gamma_h(G_{n,d})\ge
 \frac n{B_h}
 \bigl(\log B_h-2\log\log B_h-\omega\bigr).
\]
At the comparison scale $B_h\asymp\sqrt n$, this is
\[
 \gamma_h(G_{n,d})\ge
 \frac n{B_h}
 \bigl(\tfrac12\log n-2\log\log n-\omega+O(1)\bigr)
 =\Omega(\sqrt n\log n).
\]
This is an annealed lower transition, not a quenched matching theorem.  Proving the existence of dominating sets at the same coordinate remains separate.

The comparison scale is also relevant to pursuit-evasion: Meyniel's conjecture predicts that $O(\sqrt n)$ cops suffice in every connected $n$-vertex graph~\cite{LuPeng2012,ScottSudakov2011}, and it is known for several random-graph models, including random regular graphs~\cite{PralatWormald2016,PralatWormald2019}.  Our conclusion is not a cop-number lower bound.  It says that one particular static exhaustive coverage mechanism can require a logarithmic factor more than the Meyniel scale.

\subsection{Why the first moment is difficult}
For an independent random subset of density $\alpha$, a fixed tree-like radius-$h$ ball is missed with probability approximately $(1-\alpha)^{B_h}$.  Balancing subset entropy against this coupon term predicts the coordinate
\[
 \alpha B_h\approx \log B_h-2\log\log B_h.
\]
Turning that heuristic into a theorem is not a product-measure calculation.  In the configuration model, coverage events for different vertices overlap heavily, and optimizing over the selected set allows highly organized profiles.  A direct bounded-differences argument is also ineffective in the growing-radius regime: changing one pairing can alter the radius-$h$ coverage status of order $B_h$ vertices, so the natural Lipschitz constant grows on exactly the scale that must be resolved.

The key device is to label every vertex by its \emph{exact} distance to the candidate set.  Local consistency of these labels---adjacent labels differ by at most one, and every positive label has a neighbor one level lower---forces them to be genuine graph distances on any graph.  This removes the need to approximate neighborhoods by trees and converts the first moment into an exact finite-dimensional method-of-types problem.  The price is a nonconcave variational functional whose global optimizer must be identified uniformly as the number of distance levels grows.

\subsection{Related work}
The configuration or pairing model and its transfer to uniformly random simple regular graphs are standard; see Wormald's survey~\cite{Wormald1999} and Janson's simplicity theorem~\cite{Janson2009}.  Fixed-radius domination in random regular graphs has been studied algorithmically: Duckworth analyzed randomized greedy algorithms for distance-$k$ dominating sets~\cite{Duckworth2005}, while Duckworth and Wormald treated independent domination~\cite{DuckworthWormald2006}.  The broader covering viewpoint is classical in coding theory~\cite{CohenHonkalaLitsynLobstein1997}.

The closest message-passing antecedents are statistical-mechanical.  Zhao, Habibulla, and Zhou developed a cavity-method and belief-propagation treatment of minimum dominating sets~\cite{ZhaoHabibullaZhou2015}; Habibulla and Qin studied a distance-two version with distance-labeled messages~\cite{HabibullaQin2019}.  Those works are replica-symmetric and fixed-radius.  The transfer equations below are closely related in spirit; the distinctions here are the exact graph-level type count, the proof of global optimization, and the growing-radius asymptotics.

Cutler and Radcliffe used Shearer entropy to obtain universal extremal bounds for domination polynomials of regular graphs~\cite{CutlerRadcliffe2016}.  Such universal bounds cannot by themselves detect the random-covering penalty: structured regular graphs may admit efficient covering codes.  In the binomial random graph, Glebov, Liebenau, and Szab\'o proved two-point concentration at the first-moment threshold in a sufficiently dense regime~\cite{GlebovLiebenauSzabo2015}.  That result motivates, but does not supply, the corresponding quenched statement for random regular graphs.

The pursuit application comes from branch-based local coverage.  Aigner and Fromme introduced the multiple-cop game in its modern graph-theoretic form~\cite{AignerFromme1984}; Pra\l at and Wormald later combined random placement with deterministic pursuit in random graphs~\cite{PralatWormald2016,PralatWormald2019}.  The local tube certificate in~\cite[Theorem~3.3 and Section~6.4]{Neuwirth2026Tube} motivates a root-independent two-witness covering problem.  Section~\ref{subsec:tube-connection} states the exact graph-general parameter used here and carefully limits the resulting obstruction.

\subsection{Proof architecture and contributions}
The proof has four acts.

\paragraph{Act I: exact types and a compact functional.}
Exact distance labels produce a two-sided type estimate
\[
 \log \E N_{\mathbf n}
 =n\Phi_{d,h}(\mathbf n/n)+O_d(h\log n).
\]
Optimizing the local profile entropy at fixed tridiagonal edge masses reduces $\Phi_{d,h}$ exactly to a compact functional $\Fcal_{d,h}$ in $2h+1$ variables.  The lower side of the type estimate also yields the pointwise entropy anchor
$\Psi_{d,h}(\alpha)\le H(\alpha)$.

\paragraph{Act II: globality without concavity.}
The compact functional is genuinely nonconcave.  Nevertheless, entropy singularities repel every maximizing profile from the boundary.  Interior KKT points are equivalent to positive message solutions.  An exact reverse transfer reconstructs every positive stationary solution from one terminal parameter, and the associated activity is strictly increasing from $0$ to $\infty$.  Hence the grand-canonical optimizer is unique at every activity, and duality gives
\[
 \Psi_{d,h}(\alpha)=\Psi^{\mathrm{stat}}_{d,h}(\alpha),
 \qquad
 \Psi_{d,h}'(\alpha)=-\log z(\alpha).
\]

\paragraph{Act III: quantitative orbit matching.}
The reverse orbit has exact stable and free outer solutions.  They overlap over linearly many levels.  Retaining the actual transverse errors gives, for
$c=\alpha B_h/\log B_h<4/3$,
\[
 -\log z(\alpha)-\log\frac1\alpha
 =B_h(1-\alpha)^{B_h}
 \left(1+B_h^{-c/4+o(1)}\right)
 +B_h^{-c/4+o(1)}.
\]
The restriction $c<4/3$ is exactly where the additive reconstruction error remains smaller than the coupon term.

\paragraph{Act IV: one terminal conditioning cost.}
An explicit capped-free profile agrees with the free Bernoulli tree law at every nonterminal layer.  Its only entropy loss comes from conditioning the terminal lower-neighbor set to be nonempty.  That loss is exactly
\[
 p_h\Delta_d(\varepsilon_h),
 \qquad
 p_h\varepsilon_h^d=(1-\alpha)^{B_h},
 \qquad
 \Delta_d(\varepsilon)=\varepsilon^d(1+o(1)).
\]
This supplies the sharp upper bound on $H-\Psi$.  Integrating the quantitative activity law supplies the matching lower bound and the universal scaling function $1-e^{-s}$.

The principal contributions are therefore: an exact growing-radius type count with no local-tree assumption; an exact compact reduction; a global solution of a nonconcave annealed variational problem; quantitative stable/free matching; and a bounded-window theorem anchored by a self-contained feasible profile.  Quenched matching and direct two-branch asymptotics remain separate problems.

\section{Models, notation, and main results}\label{sec:main}
Fix $d\ge3$ throughout and put
\[
 b=d-1,\qquad D=\frac{b+1}{b-1}=\frac d{d-2},
\]
\begin{equation}\label{eq:Bh}
 B_h=1+d\frac{b^h-1}{b-1}=Db^h-\frac2{b-1}.
\end{equation}
For a probability vector $r=(r_1,\ldots,r_k)$, write
\[
 H(r)=-\sum_{j=1}^k r_j\log r_j,
\]
with $0\log0=0$; for a scalar $a\in[0,1]$, $H(a)$ denotes the binary entropy $H(a,1-a)$.

\begin{center}
\begin{tabular}{ll}
\toprule
Notation & Meaning\\
\midrule
$B_h$ & maximum radius-$h$ ball volume at degree $d$\\
$\gamma_h(G)$ & minimum size of a distance-$h$ dominating set\\
$\Phi_{d,h}$ & full local-profile type functional\\
$\Fcal_{d,h}$ & compact tridiagonal profile functional\\
$\Psi_{d,h}(\alpha)$ & microcanonical annealed exponent at density $\alpha$\\
$z=e^\theta$ & grand-canonical activity\\
$b=d-1$, $D=d/(d-2)$ & recurring degree constants\\
\bottomrule
\end{tabular}
\end{center}

\begin{definition}[Internally two-path $(h,2)$ domination]\label{def:two-path}
A set $S\subseteq V(G)$ is \emph{internally two-path $(h,2)$ dominating} if every vertex $v\notin S$ has two $v$--$S$ paths of length at most $h$ whose only common vertex is $v$.  In particular, the paths begin through distinct neighbors of $v$.  Let $\gamma_{h,2}^{\mathrm{int}}(G)$ denote the minimum size of such a set.
\end{definition}

Every internally two-path $(h,2)$-dominating set is distance-$h$ dominating, since either witness path alone ends in $S$ within distance $h$.  Therefore
\begin{equation}\label{eq:two-path-dominates}
 \gamma_{h,2}^{\mathrm{int}}(G)\ge \gamma_h(G).
\end{equation}

\subsection{Connection with static tube coverage}\label{subsec:tube-connection}
In the tree-ball setting of~\cite{Neuwirth2026Tube}, a length-$t$ tube with residual depth $h=R-t$ depends, after the root is allowed to vary, only on its terminal directed edge.  The resulting forward cone excludes one branch at its head.  A root-independent set meets every such directed cone exactly when, at every unselected vertex, at least two distinct branches contain a selected vertex within distance $h$; in a tree-ball these are precisely the two internally disjoint paths of Definition~\ref{def:two-path}.  Thus internally two-path domination is a graph-general strengthening of the global exhaustive tube certificate.

Combining~\eqref{eq:two-path-dominates} with the main theorem shows that any such exhaustive static certificate has size $\Omega(\sqrt n\log n)$ at the comparison scale $B_h\asymp\sqrt n$.  Meyniel's conjecture concerns the existence of a fully adaptive winning strategy with $O(\sqrt n)$ cops, so there is no contradiction: the lower bound applies only to this static exhaustive mechanism.  Partial coverage, adaptive reassignment, and epoch-based pursuit remain outside the argument.

The pairing model consists of $n$ labeled buckets of $d$ half-edges paired uniformly at random; conditioning the resulting multigraph on simplicity gives the uniformly random simple $d$-regular graph $G_{n,d}$, for $dn$ even.  Let $Z_{n,d,h}(m)$ count distance-$h$ dominating $m$-sets in the pairing model.

For a distance-$h$ dominating set $S$, assign each vertex its exact label $\operatorname{dist}(v,S)\in\{0,\ldots,h\}$.  The local consistency used below is equivalent to genuine distance: adjacent labels differ by at most one, and every positive label has a neighbor one level lower.  The descending condition gives a path to label zero of the displayed length, while the Lipschitz condition gives the reverse inequality.  No tree-neighborhood assumption is involved.

The type count developed below gives
\begin{equation}\label{eq:type-count-intro}
 \log\E Z_{n,d,h}(m)
 \le n\Psi_{d,h}(m/n)+O_d(h\log(n+1)),
\end{equation}
uniformly in $h$.  Here $\Psi_{d,h}$ is the exact annealed variational exponent.

\begin{theorem}[Bounded annealed critical window]\label{thm:bounded-critical-window}
Put $L_h=\log B_h$.  For every fixed $M<\infty$, uniformly for
\[
 C=L_h-2\log L_h+s,
 \qquad |s|\le M,
 \qquad \alpha=C/B_h,
\]
one has
\begin{equation}\label{eq:defect-bounded-window}
 H(\alpha)-\Psi_{d,h}(\alpha)
 =(1-\alpha)^{B_h}\left(1+o(1)\right).
\end{equation}
More precisely,
\begin{equation}\label{eq:defect-bounded-two-sided}
 1-B_h^{-1/7+o(1)}
 \le
 \frac{H(\alpha)-\Psi_{d,h}(\alpha)}{(1-\alpha)^{B_h}}
 \le
 1+B_h^{-(d-2)/d+o(1)}.
\end{equation}
All $o(1)$ terms are uniform over the displayed window.
\end{theorem}

\begin{corollary}[Critical scaling function]\label{cor:critical-scaling-function}
Uniformly for bounded $s$,
\begin{equation}\label{eq:critical-scaling-function}
 \boxed{
 \frac{B_h}{(\log B_h)^2}
 \Psi_{d,h}\!\left(
 \frac{\log B_h-2\log\log B_h+s}{B_h}
 \right)
 \longrightarrow 1-e^{-s}.
 }
\end{equation}
\end{corollary}

\begin{corollary}[Annealed zero and scalar expansion]\label{cor:annealed-zero-expansion}
Let
\[
 \alpha_h^{\rm ann}
 =\inf\{\alpha:\Psi_{d,h}(\alpha)\ge0\},
 \qquad
 C_h^{\rm ann}=B_h\alpha_h^{\rm ann},
\]
and let $\widehat C_{B_h}$ be the solution near
$\log B_h-2\log\log B_h$ of
\begin{equation}\label{eq:scalar-coupon-root}
 H(C/B_h)=(1-C/B_h)^{B_h}.
\end{equation}
Then
\begin{equation}\label{eq:graph-scalar-zero}
 C_h^{\rm ann}
 =\widehat C_{B_h}+O\!\left(B_h^{-1/7+o(1)}\right).
\end{equation}
Writing $L=\log B_h$ and $\ell=\log L$,
\begin{align}
 C_h^{\rm ann}
 ={}&L-2\ell
 +\frac{3\ell-1}{L}
 +\frac{5\ell^2-12\ell+3}{2L^2}\notag\\
 &+\frac{18\ell^3-81\ell^2+84\ell-17}{6L^3}
 +O\!\left(\frac{\ell^4}{L^4}\right).
 \label{eq:annealed-zero-expansion}
\end{align}
\end{corollary}

\begin{theorem}[Fixed-slack random-regular lower bound]\label{thm:random-fixed-window}
Fix $\omega>0$.  If $h=h(n)\to\infty$ and
\begin{equation}\label{eq:fixed-window-growth}
 \frac{n(\log B_h)^2}{B_h}\gg h\log n,
\end{equation}
then, with high probability,
\begin{equation}\label{eq:random-fixed-window}
 \gamma_h(G_{n,d})
 \ge
 \frac n{B_h}
 \left(\log B_h-2\log\log B_h-\omega\right).
\end{equation}
The same conclusion holds for every stronger feasible-set notion, including internally two-path $(h,2)$ domination.
\end{theorem}

The following older-form consequences retain uniform information deeper in the subcritical region and will be useful for comparison.

\begin{theorem}[Near-critical annealed negativity]\label{thm:near-critical}
Fix $\eta\in(0,1)$, put $L_h=\log B_h$, and let $W_h\to\infty$.  Uniformly for real $C$ satisfying
\begin{equation}\label{eq:C-window}
 \eta L_h\le C\le L_h-2\log L_h-W_h,
\end{equation}
one has
\begin{equation}\label{eq:near-critical-psi}
 \Psi_{d,h}(C/B_h)
 \le-\left(\frac12-o(1)\right)e^{-C}
 \qquad(h\to\infty),
\end{equation}
where the $o(1)$ may depend on $d$, $\eta$, and the prescribed sequence $W_h$, but is uniform over the displayed interval.
\end{theorem}

\begin{remark}[The coefficient $1/2$]\label{rem:half-artifact}
The coefficient $1/2$ is not predicted to be sharp.  It comes from the symmetric choice of the upper integration point in Section~\ref{sec:integration}.  The activity law and the diagnostics in Section~\ref{sec:numerics} are consistent with the normalized ratio approaching $1-O(e^{-W_h})$ deeper in the near-critical window.
\end{remark}

\begin{corollary}[Fixed-fraction slack]\label{cor:fixed-slack}
For every fixed $0<c_-\le c_+<1$, uniformly for $c\in[c_-,c_+]$,
\[
 \Psi_{d,h}\!\left(c\frac{\log B_h}{B_h}\right)
 \le-\left(\frac12-o(1)\right)B_h^{-c}.
\]
\end{corollary}

\begin{theorem}[Random-regular near-critical lower bound]\label{thm:random-near-critical}
Let $h=h(n)\to\infty$ and $W_h\to\infty$.  Define
\[
 C_h^*=L_h-2\log L_h-W_h,
\]
and suppose
\begin{equation}\label{eq:near-growth-condition}
 \liminf_{n\to\infty}\frac{C_h^*}{L_h}>0,
 \qquad
 \frac{n e^{W_h}L_h^2}{B_h}\gg h\log n.
\end{equation}
Then, with high probability,
\begin{equation}\label{eq:near-gamma-lower}
 \gamma_h(G_{n,d})\ge
 \frac n{B_h}\left(L_h-2\log L_h-W_h\right).
\end{equation}
The same conclusion holds for every parameter whose feasible sets are necessarily distance-$h$ dominating, including internally two-path $(h,2)$ domination.
\end{theorem}

\begin{corollary}[Fixed-fraction random lower bound]\label{cor:random-fixed}
For every fixed $\eps\in(0,1)$, if
\[
 \frac{n}{B_h^{1-\eps}}\gg h\log n,
\]
then, with high probability,
\[
 \gamma_h(G_{n,d})\ge
 (1-\eps)\frac{n\log B_h}{B_h}.
\]
\end{corollary}

When $B_h\asymp\sqrt n$, one has $\log B_h=\tfrac12\log n+O(1)$, and Theorem~\ref{thm:random-near-critical} gives
\[
 \gamma_h(G_{n,d})\ge
 \frac n{B_h}\left(\frac12\log n-2\log\log n-W_h+O(1)\right)
 =\Omega(\sqrt n\log n).
\]
The implicit constant in the final order statement depends on the comparison constants in $B_h\asymp\sqrt n$; no exact prefactor $\sqrt n$ is asserted.


\section{Two-sided exact type counting and the entropy anchor}\label{sec:two-sided-types}
For completeness, we record the exact local-profile count that underlies both the variational formula and the pointwise entropy bound used later.

Let $\mathcal T_{d,h}$ be the finite set of local profiles
\[
 \tau=(i,\mathbf c),\qquad
 i\in\{0,\ldots,h\},\quad
 \mathbf c=(c_0,\ldots,c_h),\quad
 \sum_jc_j=d,
\]
with
\[
 c_j=0\quad\text{if }|i-j|>1,
 \qquad
 c_{i-1}\ge1\quad\text{if }i>0.
\]
For a probability vector $\xi=(\xi_\tau)_{\tau\in\mathcal T_{d,h}}$, put
\[
 q_{ij}(\xi)
 =\frac1d\sum_{\tau=(i,\mathbf c)}\xi_\tau c_j,
\]
and call $\xi$ feasible when $q_{ij}=q_{ji}$.  Its selected density is
\[
 \alpha(\xi)=\sum_{\tau:\,i(\tau)=0}\xi_\tau.
\]
Define
\begin{equation}\label{eq:full-profile-functional}
 \Phi_{d,h}(\xi)
 =-\sum_\tau\xi_\tau\log\xi_\tau
 +\sum_\tau\xi_\tau\log\binom d{\mathbf c(\tau)}
 +\frac d2\sum_{i,j}q_{ij}(\xi)\log q_{ij}(\xi),
\end{equation}
with $0\log0=0$.  The exact variational exponent is
\[
 \Psi_{d,h}(\alpha)
 =\max\{\Phi_{d,h}(\xi):\xi\text{ feasible},\ \alpha(\xi)=\alpha\}.
\]
Section~\ref{sec:compact-reduction} proves the exact reduction from this full profile functional to the compact tridiagonal functional~\eqref{eq:compact}.

\begin{theorem}[Two-sided count for one integer type]\label{thm:two-sided-type}
There is a constant $K_d$ such that the following holds for every $h,n$ and every admissible integer profile $(n_\tau)_{\tau\in\mathcal T_{d,h}}$.  Put
\[
 \xi_\tau=\frac{n_\tau}{n},
 \qquad
 H_{ij}=\sum_{\tau=(i,\mathbf c)}n_\tau c_j=dnq_{ij}(\xi).
\]
Assume $\sum_\tau n_\tau=n$, $H_{ij}=H_{ji}$, and every $H_{ii}$ is even; these conditions define admissibility here.  Let $N_{\mathbf n}$ denote the number of vertex sets whose exact-distance local-profile counts equal $\mathbf n$ in the $d$-regular pairing model.  Then
\begin{equation}\label{eq:exact-type-count}
 \E N_{\mathbf n}
 =\frac{n!}{\prod_\tau n_\tau!}
  \prod_\tau\binom d{\mathbf c(\tau)}^{n_\tau}
  \frac{
   \displaystyle\prod_{0\le i<j\le h}H_{ij}!
   \prod_{i=0}^h(H_{ii}-1)!!
  }{(dn-1)!!},
\end{equation}
and
\begin{equation}\label{eq:two-sided-type}
 \left|
  \log\E N_{\mathbf n}-n\Phi_{d,h}(\xi)
 \right|
 \le K_d(h+1)\log(n+1).
\end{equation}
Here $(-1)!!=1$ when a diagonal mass is zero.
\end{theorem}

\begin{proof}
First assign the $n$ labeled vertices to their local-profile classes, giving $n!/\prod_\tau n_\tau!$ choices.  At a vertex of profile $(i,\mathbf c)$, assign the $d$ labeled half-edges their neighbor labels in $\binom d{\mathbf c}$ ways.  There are then $H_{ij}$ half-edges of ordered type $(i,j)$.  For $i<j$, pairing the $(i,j)$ half-edges bijectively with the $(j,i)$ half-edges gives $H_{ij}!$ choices; the $H_{ii}$ diagonal half-edges admit $(H_{ii}-1)!!$ pairings.  Division by the total number $(dn-1)!!$ of pairings proves~\eqref{eq:exact-type-count}.

Every resulting pairing has the prescribed exact distance labels.  Indeed, the allowed edge types make the label function $1$-Lipschitz, while every positive label has a neighbor one level lower.  Following a descending edge reaches label zero in exactly the displayed number of steps, and following any path from label zero cannot increase the label by more than one per edge.  Thus the label equals graph distance to its zero set.

Use, uniformly for integers $r\ge0$,
\[
 \log(r!)=r\log r-r+O(\log(r+1))
\]
and, for even $r$,
\[
 \log((r-1)!!)=\frac r2\log r-\frac r2+O(\log(r+1)).
\]
There are $O_d(h+1)$ profile classes and nonzero tridiagonal edge types.  Substitution into~\eqref{eq:exact-type-count} cancels every $\log n$, $\log d$, and linear term, leaving exactly $n\Phi_{d,h}(\xi)$ with error $O_d(h\log(n+1))$.  This proves~\eqref{eq:two-sided-type}.
\end{proof}

\begin{corollary}[Pointwise subset-entropy bound]\label{cor:Psi-le-H}
For every $d,h$ and every $\alpha\in[0,1]$,
\begin{equation}\label{eq:Psi-le-H}
 \boxed{\Psi_{d,h}(\alpha)\le H(\alpha).}
\end{equation}
\end{corollary}

\begin{proof}
First let $\xi$ be a rational feasible profile.  Passing to arbitrarily large multiples of a common denominator, and multiplying once more if necessary to make every diagonal half-edge count even, realizes $\xi$ as an integer type.  Since $N_{\mathbf n}$ counts only subsets of size $\alpha n$,
\[
 \E N_{\mathbf n}\le\binom n{\alpha n}.
\]
The lower inequality in~\eqref{eq:two-sided-type}, followed by $n\to\infty$, gives
\[
 \Phi_{d,h}(\xi)\le H(\alpha).
\]
The feasible set is a rational polytope, so rational feasible profiles are dense.  The functional $\Phi_{d,h}$ is continuous on it under the convention $0\log0=0$.  Hence the same inequality holds for every feasible profile.  Taking the maximum at fixed selected density proves~\eqref{eq:Psi-le-H}.
\end{proof}

\begin{remark}[Uniform upper type estimate]\label{rem:upper-type}
Summing~\eqref{eq:two-sided-type} over the at most $(n+1)^{O_d(h)}$ integer profiles gives
\[
 \log\E Z_{n,d,h}(m)
 \le n\Psi_{d,h}(m/n)+O_d(h\log(n+1)),
\]
uniformly as $h$ varies.  Theorem~\ref{thm:two-sided-type} also gives a matching lower estimate along every sequence realizing a prescribed rational profile.  No local-tree assumption is involved.
\end{remark}


\section{The exact compact variational problem}\label{sec:compact}
Let $q_{ij}$ be the directed-edge distribution of exact distance labels $i,j\in\{0,\dots,h\}$.  Exact distance consistency makes $q$ symmetric and tridiagonal.  Write
\[
 x_i=q_{i-1,i}\quad(1\le i\le h),
 \qquad
 \ell_i=q_{ii}\quad(0\le i\le h),
\]
with $x_0=x_{h+1}=0$.  The label masses are
\[
 p_i=\ell_i+x_i+x_{i+1},
 \qquad
 \sum_{i=0}^h p_i=1.
\]
For $i\ge1$, put
\[
 y_i=p_i-x_i=\ell_i+x_{i+1},
 \qquad
 a_i=\frac{x_i}{p_i}.
\]
Every positive-label vertex has at least one lower-label neighbor, so
\[
 a_i\ge\frac1d,
 \qquad\text{equivalently}\qquad
 d x_i\ge p_i.
\]

Let $s_d(a)$ be the maximum entropy of a distribution on nonempty subsets of $[d]$ for which each coordinate is present with marginal $a$.  Its dual formula is
\begin{equation}\label{eq:sd}
 s_d(a)=\inf_{\lambda>0}
 \left\{
 \log\bigl((1+\lambda)^d-1\bigr)-da\log\lambda
 \right\},
 \qquad \frac1d\le a\le1.
\end{equation}
The minimizing parameter is characterized by
\begin{equation}\label{eq:sd-marginal}
 a=\frac{\lambda(1+\lambda)^{d-1}}{(1+\lambda)^d-1}.
\end{equation}


\subsection{Exact local entropy reduction}\label{sec:compact-reduction}
The passage from the full profile functional to $\Fcal_{d,h}$ is exact.

\begin{lemma}[Conditional-product reduction]\label{lem:compact-reduction}
Fix a symmetric tridiagonal directed-edge distribution $q$, and write $x_i,\ell_i,p_i,y_i,a_i$ as above, with $x_{h+1}=0$.  Among all feasible local-profile laws inducing $q$, the maximum of the vertex entropy and arrangement terms in~\eqref{eq:full-profile-functional} is
\begin{align}
 &H(p_0,\ldots,p_h)
 +p_0 d\,H\!\left(\frac{\ell_0}{p_0},\frac{x_1}{p_0}\right)\notag\\
 &\quad+\sum_{i=1}^h p_i\left[
 s_d(a_i)+d(1-a_i)
 H\!\left(\frac{\ell_i}{y_i},\frac{x_{i+1}}{y_i}\right)
 \right].\label{eq:local-reduction}
\end{align}
The maximizing law is unique whenever all displayed masses are positive.  Conditional on label $i\ge1$, the set of lower-label half-edges has the entropy-maximizing tilted nonempty-subset law with coordinate marginal $a_i$; conditional on that set, every remaining half-edge independently receives label $i$ or $i+1$ with probabilities $\ell_i/y_i$ and $x_{i+1}/y_i$.
\end{lemma}

\begin{proof}
First expose the central label.  Its entropy is $H(p_0,\ldots,p_h)$.  Conditional on the central label $i$, choosing a local count vector and then assigning the $d$ labeled half-edges contributes exactly the entropy of the induced law on words of length $d$ over the allowed neighboring labels.

For $i=0$, only labels $0$ and $1$ are allowed and their coordinate marginals are $\ell_0/p_0$ and $x_1/p_0$.  Subadditivity of entropy is sharp only for independent coordinates, giving the first term of~\eqref{eq:local-reduction}.

Fix $i\ge1$.  Mark the coordinates whose neighbor label is $i-1$.  Their random subset is nonempty and has common coordinate marginal
\[
 a_i=\frac{q_{i,i-1}}{p_i}=\frac{x_i}{p_i}.
\]
By definition, its entropy is at most $s_d(a_i)$, with equality for the exponential tilt
\[
 \Pr(A)\propto \lambda^{|A|},\qquad \varnothing\ne A\subseteq[d],
\]
where $\lambda$ satisfies~\eqref{eq:sd-marginal}; this also proves the dual formula~\eqref{eq:sd}.  Given the lower-neighbor set, the remaining $d-|A|$ coordinates must split between labels $i$ and $i+1$.  Conditional entropy is maximized by independent splitting with probabilities $\ell_i/y_i$ and $x_{i+1}/y_i$.  Its expectation is
\[
 d(1-a_i)H\!\left(\frac{\ell_i}{y_i},\frac{x_{i+1}}{y_i}\right).
\]
The chain rule for entropy proves~\eqref{eq:local-reduction}; strict entropy concavity gives uniqueness in the positive interior.
\end{proof}

\begin{proposition}[Exact compact functional]\label{prop:exact-compact}
For every feasible tridiagonal $q$, maximizing~\eqref{eq:full-profile-functional} over local-profile laws inducing $q$ gives exactly~\eqref{eq:compact}.  Consequently the compact and full variational values agree at every selected density.
\end{proposition}

\begin{proof}
Insert~\eqref{eq:local-reduction} into~\eqref{eq:full-profile-functional}.  The edge term is
\[
 \frac d2\left(\sum_{i=0}^h\ell_i\log\ell_i+2\sum_{i=1}^h x_i\log x_i\right).
\]
Expanding the two categorical entropies in~\eqref{eq:local-reduction}, the $x_i\log x_i$ terms cancel against the off-diagonal edge terms.  The remaining $p_i$, $y_i$, and $\ell_i$ terms collect to~\eqref{eq:compact}, including the root contribution $(d-1)p_0\log p_0$.  Since every full feasible profile induces a feasible $q$ and Lemma~\ref{lem:compact-reduction} constructs a maximizing profile for every feasible $q$, the variational values coincide.
\end{proof}

Define the feasible polytope
\[
 \Pcal_{d,h}
 =\left\{(x,\ell):
 x_i,\ell_i\ge0,
 \ \sum_i p_i=1,
 \ d x_i\ge p_i\ (1\le i\le h)
 \right\}.
\]
With the convention $0\log0=0$, the exact microcanonical exponent at a profile is
\begin{equation}\label{eq:compact}
\begin{aligned}
 \Fcal_{d,h}(x,\ell)
 ={}&(d-1)p_0\log p_0\\
 &+\sum_{i=1}^h
 \left[
 -p_i\log p_i+p_i s_d(a_i)+d y_i\log y_i
 \right]
 -\frac d2\sum_{i=0}^h\ell_i\log\ell_i.
\end{aligned}
\end{equation}
For activity $z=e^\theta>0$, the grand-canonical functional is
\begin{equation}\label{eq:grand-functional}
 \Fcal_{d,h}^{(z)}(x,\ell)
 =\Fcal_{d,h}(x,\ell)+p_0\log z.
\end{equation}
The exact variational values are
\[
 \Psi_{d,h}(\alpha)
 =\max_{\substack{(x,\ell)\in\Pcal_{d,h}\\p_0=\alpha}}
 \Fcal_{d,h}(x,\ell),
\]
\[
 \phi_{d,h}(z)
 =\max_{(x,\ell)\in\Pcal_{d,h}}
 \Fcal_{d,h}^{(z)}(x,\ell)
 =\max_\alpha\{\Psi_{d,h}(\alpha)+\alpha\log z\}.
\]
The polytope is compact, so all maxima exist.

\subsection{The feasible density interval}
The local capacities imply the Moore lower bound.  First,
\[
 p_1\le d x_1\le d p_0.
\]
For $i\ge2$, symmetry and the preceding descent edge give
\[
 x_i\le p_{i-1}-x_{i-1}
 \le\frac{d-1}{d}p_{i-1},
\]
so
\[
 p_i\le d x_i\le(d-1)p_{i-1}.
\]
Therefore
\[
 1=\sum_{i=0}^h p_i\le B_h p_0,
 \qquad p_0\ge\frac1{B_h}.
\]
Equality is feasible in the type polytope: take
\[
 p_0=\frac1{B_h},
 \qquad
 p_i=\frac{d(d-1)^{i-1}}{B_h},
 \qquad
 x_i=\frac{(d-1)^{i-1}}{B_h},
\]
with $\ell_i=0$ for $i<h$ and $\ell_h=(d-1)^h/B_h$.  The upper endpoint $p_0=1$ is also feasible.  Hence the density projection of $\Pcal_{d,h}$ is exactly
\[
 \left[\frac1{B_h},1\right].
\]


\section{Boundary repulsion and full interiority}
The proof that all grand-canonical maximizers are interior rests on the endpoint behavior of $s_d$.

\begin{lemma}[Endpoint expansions]\label{lem:sd-endpoints}
As $\delta\downarrow0$,
\begin{align}
 s_d\left(\frac1d+\delta\right)
 &=\log d+d\delta\log\frac1\delta+O_d(\delta),
 \label{eq:sd-lower}\\
 s_d(1-\delta)
 &=d\delta\log\frac1\delta+O_d(\delta).
 \label{eq:sd-upper}
\end{align}
\end{lemma}

\begin{proof}
The envelope theorem applied to~\eqref{eq:sd} gives
\[
 s_d'(a)=-d\log\lambda(a).
\]
Expanding~\eqref{eq:sd-marginal} at $\lambda=0$ gives
\[
 \lambda(a)=\frac{2d}{d-1}
 \left(a-\frac1d\right)
 +O_d\left(\left(a-\frac1d\right)^2\right).
\]
At the other endpoint, $\lambda(a)=(1-a)^{-1}(1+O_d(1-a))$.  Integrating $s_d'$ from the known endpoint values $s_d(1/d)=\log d$ and $s_d(1)=0$ gives the result.
\end{proof}

If $p_i=0$ at one level, then $x_{i+1}=0$, and the capacity inequality forces $p_{i+1}=0$.  Thus every feasible profile has an effective horizon
\[
 k=\max\{i:p_i>0\},
\]
and its positive support is the prefix $\{0,\dots,k\}$.

\begin{lemma}[Relative boundary repulsion]\label{lem:relative-boundary}
Fix a finite activity $z>0$.  A maximizer of $\Fcal_{d,h}^{(z)}$ cannot lie on any proper local face within its effective horizon.  More precisely, if its effective horizon is $k$, then
\[
 \ell_i>0\quad(0\le i\le k),
 \qquad
 a_i>\frac1d\quad(1\le i\le k).
\]
In particular, $a_k<1$.
\end{lemma}

\begin{proof}
For each $k\ge1$, a strictly interior feasible profile exists.  Choose
\[
 \frac1d<\chi<\frac12,
 \qquad
 0<r<\chi^{-1}-1,
\]
put $p_i$ proportional to $r^i$, and set $x_i=\chi p_i$.  Then
\[
 \ell_0=p_0-x_1>0,
 \quad
 \ell_i=p_i-x_i-x_{i+1}>0\ (1\le i<k),
 \quad
 \ell_k=p_k-x_k>0.
\]

Mix a proposed boundary maximizer with such an interior profile.  If $a_i=1/d$, equation~\eqref{eq:sd-lower} contributes a strictly positive multiple of $t\log(1/t)$.  If $\ell_i=0$, the term
\[
 -\frac d2\ell_i\log\ell_i
\]
contributes a strictly positive multiple of $t\log(1/t)$.  All terms that remain away from their endpoints change by only $O(t)$.

The only apparent negative singularity occurs if $a_k=1$, equivalently $y_k=\ell_k=0$.  In that case, with $y_k(t)=\beta t+O(t^2)$ for some $\beta>0$,
\[
 p_k(t)s_d\left(1-\frac{y_k(t)}{p_k(t)}\right)
 +d y_k(t)\log y_k(t)=O(t)
\]
by~\eqref{eq:sd-upper}; the two logarithmic singularities cancel.  The remaining diagonal-edge term contributes
\[
 -\frac d2y_k(t)\log y_k(t)
 =\frac d2\beta\,t\log\frac1t+O(t),
\]
which is strictly positive.  If several faces are active simultaneously, all leading $t\log(1/t)$ coefficients are nonnegative and at least one is positive.  Hence every proper relative boundary point admits an improving inward direction.
\end{proof}

The horizon itself is also repelling.

\begin{lemma}[Horizon extension]\label{lem:horizon-extension}
A maximizer of $\Fcal_{d,h}^{(z)}$ cannot have effective horizon $k<h$.
\end{lemma}

\begin{proof}
By Lemma~\ref{lem:relative-boundary}, a maximizer is strictly interior on its effective prefix.  Choose
\[
 a_*\in\left(\max\left\{\frac1d,1-\frac2d\right\},1\right).
\]
For small $\tau>0$, introduce a new level of mass $p_{k+1}=\tau$ with
\[
 x_{k+1}=a_*\tau,
 \qquad
 y_{k+1}=\ell_{k+1}=(1-a_*)\tau.
\]
Keep $p_k$ fixed by reducing $\ell_k$ by $a_*\tau$, and preserve total mass by reducing $\ell_0$ by $\tau$.  All old variables remain feasible for sufficiently small $\tau$.  When $k=0$, take instead
\[
 p_0=1-\tau,
 \quad
 x_1=a_*\tau,
 \quad
 \ell_0=1-(1+a_*)\tau.
\]

The new level contributes
\[
 \left[1-\frac d2(1-a_*)\right]
 \tau\log\frac1\tau+O(\tau).
\]
The coefficient is positive by the choice of $a_*$.  All changes to previously positive coordinates are $O(\tau)$.  Thus the extension strictly increases the grand functional.
\end{proof}

\begin{theorem}[Full interiority]\label{thm:full-interiority}
For every $d\ge3$, $h\ge1$, and $z>0$, every maximizer of the exact grand-canonical functional~\eqref{eq:grand-functional} has full effective horizon $h$ and satisfies
\[
 x_i>0,
 \quad
 \ell_i>0,
 \quad
 a_i>\frac1d
\]
for all applicable indices.
\end{theorem}

\begin{proof}
Combine Lemmas~\ref{lem:relative-boundary} and~\ref{lem:horizon-extension}.
\end{proof}



\section{Stationary messages and exact KKT reconstruction}\label{sec:stationarity}
Recall $b=d-1$.  For $1\le i\le h$, let $A_i$ be the cavity weight when the recipient edge already supplies a lower-label neighbor, and let $B_i$ be the weight when it does not.  For label zero use $B_0$, and set $A_{h+1}=0$.  Define
\[
 S_0=B_0+A_1,
 \qquad
 S_i=B_{i-1}+B_i+A_{i+1}\quad(1\le i\le h),
\]
where the terminal convention makes $S_h=B_{h-1}+B_h$.  The positive stationary system is
\begin{align}
 \kappa B_0&=zS_0^b,\label{eq:msg-root}\\
 \kappa A_i&=S_i^b,&&1\le i\le h,\label{eq:msg-A}\\
 \kappa B_i&=S_i^b-(S_i-B_{i-1})^b,&&1\le i\le h.\label{eq:msg-B}
\end{align}
The following subsection derives this system directly from the compact exact functional and proves the converse reconstruction.


\subsection{Full KKT-to-message correspondence}\label{sec:kkt}
We now derive the stationary message system directly from the compact KKT equations and prove the converse reconstruction.

Let $\lambda_i$ be the minimizer in~\eqref{eq:sd}.  Then
\begin{equation}\label{eq:sd-identities}
 s_d'(a_i)=-d\log\lambda_i,
 \qquad
 s_d(a_i)-a_i s_d'(a_i)=\log((1+\lambda_i)^d-1).
\end{equation}
At an interior grand-canonical KKT point, let $\mu$ be the multiplier for
\[
 \sum_{i=0}^h\ell_i+2\sum_{i=1}^h x_i=1.
\]
Since $\ell_i$ enters one layer mass and $x_i$ enters two, the KKT equations are
\begin{equation}\label{eq:kkt-weight}
 \frac{\partial\Fcal^{(z)}}{\partial\ell_i}=\mu,
 \qquad
 \frac{\partial\Fcal^{(z)}}{\partial x_i}=2\mu.
\end{equation}

The analytic derivatives used below are, at the root,
\[
 g_0=(d-1)(\log p_0+1)+\log z-\frac d2(\log\ell_0+1),
\]
and, for $i\ge1$,
\[
 g_i=-\log p_i-1+s_d(a_i)-a_i s_d'(a_i)
 +d(\log y_i+1)-\frac d2(\log\ell_i+1).
\]
Thus $\partial\Fcal^{(z)}/\partial\ell_i=g_i$.  The derivative with respect to $x_i$ is the sum of the layer contribution immediately to its left and the $x_i$-derivative of the layer-$i$ contribution; explicitly, for $i\ge2$,
\[
 \frac{\partial\Fcal^{(z)}}{\partial x_i}
 =g_{i-1}+\frac d2(\log\ell_{i-1}+1)
 +g_i+s_d'(a_i)-d(\log y_i+1)+\frac d2(\log\ell_i+1),
\]
with the same formula at $i=1$ after replacing the left layer expression by its root analogue.

Define, up to a common positive scale,
\begin{equation}\label{eq:belief-messages}
 B_i=\sqrt{\ell_i},
 \qquad
 A_i=\frac{x_i}{B_{i-1}}
 \quad(1\le i\le h),
 \qquad
 A_{h+1}=0.
\end{equation}
Then
\[
 y_i=B_i(B_i+A_{i+1}).
\]

\begin{lemma}[Edge KKT identifies the local tilt]\label{lem:kkt-lambda}
For every $1\le i\le h$,
\begin{equation}\label{eq:lambda-message}
 \lambda_i
 =\frac{B_{i-1}}{B_i+A_{i+1}}.
\end{equation}
\end{lemma}

\begin{proof}
For $i\ge2$, the derivative of $\Fcal$ with respect to $x_i$ is the sum of the $p_{i-1}$ and $p_i$ contributions.  Subtract the two diagonal equations in~\eqref{eq:kkt-weight}, use~\eqref{eq:sd-identities}, and cancel the common constants.  The result is
\[
 \frac d2\log\ell_{i-1}
 -d\log\left(\lambda_i\frac{y_i}{\sqrt{\ell_i}}\right)=0.
\]
Thus
\[
 \lambda_i y_i=\sqrt{\ell_{i-1}\ell_i},
\]
which is~\eqref{eq:lambda-message}.  The root contribution at $i=1$ has the same algebra, with the activity term cancelling through the root diagonal equation.
\end{proof}

\begin{proposition}[Compact KKT equals two-message stationarity]\label{prop:kkt-message}
Every strict interior KKT point determines positive messages satisfying~\eqref{eq:msg-root}--\eqref{eq:msg-B}.  Conversely, every positive message solution reconstructs a strict interior KKT point by
\[
 q_{ii}=\frac{B_i^2}{Z_e},
 \qquad
 q_{i-1,i}=\frac{B_{i-1}A_i}{Z_e}.
\]
\end{proposition}

\begin{proof}
Put
\[
 T_i=B_i+A_{i+1},
 \qquad
 S_i=B_{i-1}+T_i.
\]
By Lemma~\ref{lem:kkt-lambda}, $\lambda_i=B_{i-1}/T_i$.  Exponentiating the diagonal KKT equation and using~\eqref{eq:sd-identities} shows that one common constant $\kappa>0$ satisfies
\[
 \kappa
 =\frac{((1+\lambda_i)^d-1)y_i^d}{p_iB_i^d}
 =\frac{S_i^d-T_i^d}{p_i}
\]
for every $i\ge1$.  Hence
\[
 p_i=\frac{S_i^d-T_i^d}{\kappa}.
\]
The descent marginal formula gives
\[
 x_i=a_ip_i
 =\frac{B_{i-1}S_i^b}{\kappa}.
\]
Since $x_i=B_{i-1}A_i$, this is~\eqref{eq:msg-A}.  Subtracting $x_i$ from $p_i$ gives
\[
 y_i=\frac{T_i(S_i^b-T_i^b)}{\kappa}.
\]
Since $y_i=B_iT_i$, this is~\eqref{eq:msg-B}.  At the root, the diagonal KKT equation gives
\[
 \kappa=\frac{zp_0^b}{B_0^d}.
\]
Because $p_0=B_0(B_0+A_1)=B_0S_0$, this is~\eqref{eq:msg-root}.

Conversely, let positive messages satisfy the stationarity equations.  Put $\sigma=Z_e^{-1/2}$ and introduce normalized messages
\[
 \widehat A_i=\sigma A_i,
 \qquad
 \widehat B_i=\sigma B_i,
 \qquad
 \widehat\kappa=\sigma^{b-1}\kappa.
\]
Then
\[
 \widehat\kappa\widehat A_i=\widehat S_i^b,
 \qquad
 \widehat\kappa\widehat B_i
 =\widehat S_i^b-(\widehat S_i-\widehat B_{i-1})^b,
 \qquad
 \widehat\kappa\widehat B_0=z\widehat S_0^b.
\]
The reconstructed beliefs are simply
\[
 \ell_i=\widehat B_i^2,
 \qquad
 x_i=\widehat B_{i-1}\widehat A_i.
\]
They have total mass one by the definition of $Z_e$.  Writing
\[
 \widehat T_i=\widehat B_i+\widehat A_{i+1},
 \qquad
 \lambda_i=\frac{\widehat B_{i-1}}{\widehat T_i},
\]
the message equations give the exact row identities
\[
 p_i=\frac{\widehat S_i^d-\widehat T_i^d}{\widehat\kappa},
 \qquad
 x_i=\frac{\widehat B_{i-1}\widehat S_i^b}{\widehat\kappa},
 \qquad
 y_i=\frac{\widehat T_i(\widehat S_i^b-\widehat T_i^b)}{\widehat\kappa}.
\]
Consequently $a_i=x_i/p_i$ is exactly the tilted conditioned-binomial marginal with parameter $\lambda_i$, and
\[
 \lambda_i y_i
 =\widehat B_{i-1}\widehat B_i
 =\sqrt{\ell_{i-1}\ell_i}.
\]
Substitution into the diagonal derivative gives, for every $i\ge1$,
\[
 g_i=\log\widehat\kappa+\frac{d-2}{2}.
\]
The root equation gives the same value for $g_0$.  Finally, the displayed edge identity and $s_d'(a_i)=-d\log\lambda_i$ make the difference between the $x_i$ derivative and $g_{i-1}+g_i$ vanish exactly.  Thus, with
\[
 \mu=\log\widehat\kappa+\frac{d-2}{2},
\]
one has
\[
 \frac{\partial\Fcal^{(z)}}{\partial\ell_i}=\mu,
 \qquad
 \frac{\partial\Fcal^{(z)}}{\partial x_i}=2\mu,
\]
which is~\eqref{eq:kkt-weight}.  Positivity makes the point strict interior, and common message scaling cancels from the beliefs.
\end{proof}



\subsection{One-dimensional positive stationary locus}
The equations are homogeneous: common message scaling changes $\kappa$ but not $z$ or the reconstructed profile.  In the gauge $\kappa=1$, choosing $B_h>0$ determines all preceding messages uniquely by
\begin{equation}\label{eq:message-backward}
 A_i=B_i+(B_i+A_{i+1})^b,
 \qquad
 B_{i-1}=A_i^{1/b}-B_i-A_{i+1},
\end{equation}
for $i=h,h-1,\ldots,1$.  Positivity is automatic because
\[
 B_{i-1}=\left((B_i+A_{i+1})^b+B_i\right)^{1/b}-(B_i+A_{i+1})>0.
\]
Thus the positive stationary locus is one dimensional before the activity is imposed.


\section{Exact reverse transfer and monotone activity shooting}
Normalize $A_1=1$ and put
\[
 \rho_i=\frac{B_i}{A_i},
 \qquad
 u_i=\frac{A_{i+1}}{A_i},
 \qquad
 v_i=\rho_i+u_i.
\]
At the terminal level, $u_h=0$, so
\[
 \rho_h=v_h=:s\in(0,1).
\]
Let
\[
 \Dcal=\{(\rho,v):0<\rho\le v\le1\}.
\]

\begin{proposition}[Exact reverse map]\label{prop:reverse-map}
Given a next-level state $(\rho',v')\in\Dcal$, define
\[
 w=(1-\rho')^{1/b},
 \qquad
 e=1-w,
\]
\[
 R=v'\frac ew,
 \qquad
 M=\left(e+\frac w{v'}\right)^b.
\]
Then its unique positive predecessor is
\begin{equation}\label{eq:reverse-map}
 \boxed{
 \rho=\frac{R}{R+M},
 \qquad
 v=\frac{1+R}{R+M}.
 }
\end{equation}
The remaining coordinates are
\begin{equation}\label{eq:reverse-out}
 u=\frac1{R+M},\qquad q=\frac{M-1}{R+M}.
\end{equation}
The map sends $\Dcal$ into itself.
\end{proposition}

\begin{proof}
The forward transfer identities imply
\[
 \frac\rho u=v'\frac ew=R.
\]
The next lower-neighbor fraction also gives
\[
 e=\frac{\rho(1-\rho)^{1/b}}
 {u^{1/b}(u+\rho)}.
\]
Substituting $\rho=Ru$ and solving yields
\[
 u=\frac1{R+M},
\]
which gives~\eqref{eq:reverse-map}.  Uniqueness follows from the algebraic solution.  Since
\[
 e+\frac w{v'}\ge e+w=1,
\]
we have $M\ge1$, and therefore
\[
 0<\rho<v\le1.
\]
\end{proof}

\begin{proposition}[Order preservation]\label{prop:order-preserving}
The reverse map is coordinatewise nondecreasing on $\Dcal$, and its $\rho$ coordinate is strictly increasing whenever either input coordinate increases.  Along the diagonal terminal family $(s,s)$, every earlier $\rho_i$ is strictly increasing in $s$, and every $v_i$ is nondecreasing.
\end{proposition}

\begin{proof}
The quantity $R=v'e/w$ is strictly increasing in both $\rho'$ and $v'$.  Put
\[
 N=e+\frac w{v'},
 \qquad M=N^b.
\]
As $v'$ increases, $N$ strictly decreases.  Since $w'(\rho')<0$ and $e'=-w'$, one has
\[
 \frac{\partial N}{\partial\rho'}
 =w'(\rho')\left(\frac1{v'}-1\right)\le0.
\]
Thus $M$ is nonincreasing in both inputs.  The function
\[
 \rho=\frac R{R+M}
\]
increases with $R$ and decreases with $M$.  Likewise
\[
 v=\frac{1+R}{R+M}
\]
increases with $R$ when $M\ge1$ and decreases with $M$.  The asserted monotonicity follows, and strictness of $\rho$ propagates under iteration.
\end{proof}

At the root, write
\[
 m=\frac{1-(1-\rho_1)^{1/b}}{(1-\rho_1)^{1/b}}.
\]
Then
\[
 r_0:=\frac{B_0}{A_1}=v_1m,
 \qquad
 \kappa=\bigl(v_1(1+m)\bigr)^b,
\]
and the root equation gives
\begin{equation}\label{eq:z-root-reverse}
 \boxed{
 z=v_1^{b+1}m
 \left(\frac{1+m}{1+v_1m}\right)^b.
 }
\end{equation}

\begin{lemma}[Strict activity monotonicity]\label{lem:z-monotone}
The root activity in~\eqref{eq:z-root-reverse} is strictly increasing in both $v_1$ and $m$.  Consequently, the terminal-to-activity map
\[
 s\longmapsto z_h(s)
\]
is strictly increasing on $(0,1)$.
\end{lemma}

\begin{proof}
Direct differentiation gives
\[
 \frac{\partial}{\partial v}\log z
 =\frac{b+1+vm}{v(1+vm)}>0,
\]
\[
 \frac{\partial}{\partial m}\log z
 =\frac1m+
 \frac{b(1-v)}{(1+m)(1+vm)}>0.
\]
The quantity $m$ is strictly increasing in $\rho_1$, so Proposition~\ref{prop:order-preserving} completes the proof.
\end{proof}

\begin{lemma}[Endpoint activities]\label{lem:z-endpoints}
For fixed $d,h$,
\[
 \lim_{s\downarrow0}z_h(s)=0,
 \qquad
 \lim_{s\uparrow1}z_h(s)=\infty.
\]
\end{lemma}

\begin{proof}
For terminal $(\rho',v')=(s,s)$ with $s\downarrow0$, one reverse step has $R=O(s^2)$ and $M=\Theta(s^{-b})$, so the predecessor tends to $(0,0)$.  The reverse map is continuous at every interior state and has the displayed boundary limit; induction over the fixed number $h-1$ of reverse steps therefore sends every earlier state to $(0,0)$.  Equation~\eqref{eq:z-root-reverse} then gives $z\to0$.

As $s\uparrow1$, one has $w\to0$, $R\to\infty$, and $M\to1$, so one reverse step tends to $(1,1)$.  The same finite-step induction sends every earlier state to $(1,1)$.  Hence $m\to\infty$, $v_1\to1$, and~\eqref{eq:z-root-reverse} gives $z\to\infty$.  Continuity of the reverse map and of~\eqref{eq:z-root-reverse} also makes $s\mapsto z_h(s)$ continuous.
\end{proof}

\begin{theorem}[Unique positive stationary point]\label{thm:unique-stationary}
For every $d\ge3$, $h\ge1$, and $z>0$, the exact grand-canonical stationarity equations have exactly one positive solution up to common message scaling.
\end{theorem}

\begin{proof}
Every positive solution has one terminal parameter $s=\rho_h\in(0,1)$ and is uniquely reconstructed by Proposition~\ref{prop:reverse-map}.  Lemmas~\ref{lem:z-monotone} and~\ref{lem:z-endpoints} show that $s\mapsto z_h(s)$ is a strictly increasing bijection from $(0,1)$ to $(0,\infty)$.
\end{proof}


For a positive stationary message profile whose reconstructed selected density is $\alpha$, define
\[
 \Psi^{\mathrm{stat}}_{d,h}(\alpha)
 :=\Fcal_{d,h}(x,\ell),
\]
where $(x,\ell)$ is its normalized compact profile.  Theorem~\ref{thm:unique-stationary} shows that this value is single-valued along the positive reverse-transfer branch.

\section{Exact grand- and microcanonical globality}
Although $\Fcal_{d,h}$ is nonconcave, the preceding interiority and uniqueness results determine its global optimizer.

\begin{theorem}[Grand-canonical globality]\label{thm:grand-globality}
For every $d\ge3$, $h\ge1$, and $z>0$, the exact grand-canonical functional $\Fcal_{d,h}^{(z)}$ has a unique global maximizer.  It is the positive stationary profile corresponding to the unique terminal parameter $s$ satisfying $z_h(s)=z$.
\end{theorem}

\begin{proof}
A maximizer exists by compactness.  Theorem~\ref{thm:full-interiority} places every maximizer in the strict interior, so every maximizer satisfies the interior stationarity equations.  Theorem~\ref{thm:unique-stationary} gives only one such point.
\end{proof}

The microcanonical problem follows by duality.

\begin{theorem}[Microcanonical globality]\label{thm:micro-globality}
For every density
\[
 \frac1{B_h}<\alpha<1,
\]
the exact microcanonical functional has a unique global maximizer, and it is the positive stationary profile on the reverse-transfer branch with selected density $\alpha$.  Equivalently,
\[
 \boxed{
 \Psi_{d,h}(\alpha)=\Psi^{\mathrm{stat}}_{d,h}(\alpha)
 }
\]
throughout the interior feasible interval.
\end{theorem}

\begin{proof}
Write $\theta=\log z$ and
\[
 \phi_{d,h}(\theta)
 =\max_{(x,\ell)\in\Pcal_{d,h}}
 \{\Fcal_{d,h}(x,\ell)+\theta p_0\}.
\]
By Theorem~\ref{thm:grand-globality}, the maximizer is unique for every $\theta$.  Danskin's theorem~\cite{RockafellarWets1998} therefore gives
\[
 \phi_{d,h}'(\theta)=\alpha(\theta),
\]
the density of that maximizer.  The derivative of a differentiable convex function is continuous here (equivalently, one may use continuity of the unique optimizer).  It is also strictly increasing.  Indeed, if $\theta_1<\theta_2$ had the same maximizing density, the two optimality inequalities would force each optimizer to maximize at both activities.  Grand-canonical uniqueness would make the profiles equal, but the root equation assigns one activity to an interior stationary profile, a contradiction.  To identify the endpoint limits, let $\theta_k\to-\infty$ and pass, by compactness, to a convergent subsequence of optimizers.  Comparing with a feasible profile of density $1/B_h$ shows that any limit must minimize $p_0$ over the polytope, hence has density $1/B_h$.  Similarly, along $\theta_k\to\infty$, comparison with the all-selected profile forces every subsequential limit to have density $1$.  Therefore
\[
 \alpha(\theta)\longrightarrow\frac1{B_h}
 \quad(\theta\to-\infty),
 \qquad
 \alpha(\theta)\longrightarrow1
 \quad(\theta\to\infty).
\]
Hence every interior density is attained.

Fix $\alpha$ and choose $\theta$ with $\alpha(\theta)=\alpha$.  For every profile $P$ of density $\alpha$,
\[
 \Fcal(P)+\theta\alpha
 \le
 \Fcal(P_\theta)+\theta\alpha,
\]
so $P_\theta$ is the microcanonical maximizer.  Uniqueness follows from grand-canonical uniqueness.
\end{proof}

\begin{corollary}[Concavity and the envelope identity]\label{cor:envelope}
The exact value function $\Psi_{d,h}$ is strictly concave on $(1/B_h,1)$ and differentiable there.  If $z(\alpha)$ is the activity exposing density $\alpha$, then
\[
 \boxed{
 \Psi_{d,h}'(\alpha)=-\log z(\alpha).
 }
\]
\end{corollary}

\begin{proof}
This is the standard differentiable Legendre correspondence produced by the unique maximizers in Theorems~\ref{thm:grand-globality} and~\ref{thm:micro-globality}.  Strict monotonicity of $\alpha(\theta)$ gives strict concavity.
\end{proof}

For later use, we collect the two conclusions as
\begin{equation}\label{eq:globality}
 \Psi_{d,h}(\alpha)=\Psi^{\rm stat}_{d,h}(\alpha),
 \qquad
 \Psi'_{d,h}(\alpha)=-\log z(\alpha).
\end{equation}

Explicit positive Hessian directions at stratified profiles do not contradict these theorems: they occur at nonstationary stratified profiles.  The compact functional is genuinely nonconcave, but no competing stationary maximum exists.


\section{Corrected stationary telescoping and root formulas}\label{sec:telescoping}
Define
\begin{align}
 Z_v&=zS_0^d+
 \sum_{i=1}^h
 \left[S_i^d-(S_i-B_{i-1})^d\right],
 \label{eq:Zv}\\
 Z_e&=\sum_{i=0}^hB_i^2
 +2\sum_{i=0}^{h-1}B_iA_{i+1}.
 \label{eq:Ze}
\end{align}
The stationary vertex and edge normalizers satisfy the following exact telescoping identity.

\begin{proposition}[Corrected telescoping identity]\label{prop:correct-telescope}
Every positive stationary solution satisfies
\[
 \boxed{Z_v=\kappa Z_e.}
\]
\end{proposition}

\begin{proof}
For $1\le i\le h$, put
\[
 V_i=S_i^d-(S_i-B_{i-1})^d.
\]
Using~\eqref{eq:msg-A}--\eqref{eq:msg-B},
\begin{align*}
 V_i
 &=\kappa\left(
 B_{i-1}A_i+B_i^2+B_iA_{i+1}
 \right).
\end{align*}
The root contribution is
\[
 zS_0^d=\kappa(B_0^2+B_0A_1).
\]
After summation, every square $B_i^2$ appears once and every cross term $B_iA_{i+1}$ appears twice: once from each adjacent vertex contribution, with the root supplying the first copy of $B_0A_1$.  This is exactly $\kappa Z_e$.
\end{proof}

Normalize $A_1=1$ and write $r=B_0$.  Proposition~\ref{prop:correct-telescope} gives
\begin{equation}\label{eq:alpha-root-correct}
 \alpha=\frac{r(1+r)}{Z_e}.
\end{equation}
The pressure is
\[
 \phi=\log Z_v-\frac d2\log Z_e,
\]
and the root equation is
\[
 \log z=\log\kappa+\log r-(d-1)\log(1+r).
\]
Eliminating $Z_e$ gives the corrected root-only formulas
\begin{equation}\label{eq:root-pressure-correct}
 \boxed{
 \phi
 =\log z
 +\frac d2\log\frac{1+r}{r}
 +\frac{d-2}{2}\log\alpha,
 }
\end{equation}
\begin{equation}\label{eq:root-micro-correct}
\boxed{
\begin{aligned}
 \Psi={}&
 (1-\alpha)\log\kappa
 -\left(\frac{d-2}{2}+\alpha\right)\log r
 +\frac{d-2}{2}\log\alpha\\
 &+\left((d-1)\alpha-\frac{d-2}{2}\right)\log(1+r).
\end{aligned}
}
\end{equation}
The corrected formulas are used throughout this paper.  Calculations made directly from $Z_v$ and $Z_e$ are unaffected, provided the power differences are evaluated stably as discussed in Section~\ref{sec:numerics}.

For later diagnostics, the same telescope gives a cancellation-free identity for the entropy defect.  Put
\[
 u=u_1=\frac{A_2}{A_1},
 \qquad c_0=\frac{d-2}{2},
\]
and define
\[
 E_{\rm root}:=\log\kappa-b\log u,
 \qquad
 N_e:=\log Z_e-(D+1)\log u.
\]

\begin{proposition}[Exact entropy-defect identity]\label{prop:defect-identity}
Every positive stationary solution of selected density $\alpha$ satisfies
\begin{equation}\label{eq:defect-identity}
 \boxed{
 H(\alpha)-\Psi_{d,h}(\alpha)
 =-(1-\alpha)\log(1-\alpha)
 -E_{\rm root}+c_0N_e
 +\alpha\log\frac z\alpha.
 }
\end{equation}
\end{proposition}

\begin{proof}
Proposition~\ref{prop:correct-telescope} gives
\[
 \phi=\log\kappa-c_0\log Z_e.
\]
Since $c_0(D+1)=b$, this is
$\phi=E_{\rm root}-c_0N_e$.  Now use
$\Psi=\phi-\alpha\log z$ and expand
$H(\alpha)=-\alpha\log\alpha-(1-\alpha)\log(1-\alpha)$.
\end{proof}



\section{Terminal logarithmic coordinate and exact outer orbits}\label{sec:terminal}
Retain the normalization $A_1=1$ and the reverse-transfer coordinates
\[
 \rho_i=\frac{B_i}{A_i},\qquad
 u_i=\frac{A_{i+1}}{A_i},\qquad
 v_i=\rho_i+u_i,\qquad
 q_i=1-v_i.
\]
At the terminal wall $u_h=0$.  We write
\begin{equation}\label{eq:terminal-t}
 \rho_h=v_h=s=1-e^{-t},\qquad q_h=e^{-t}.
\end{equation}
The exact reverse dynamics are those of Proposition~\ref{prop:reverse-map}.

Two boundary orbits are exact.  On the stable boundary $\rho=0$, the forward map is
$u\mapsto u^{1/b}$ and the reverse map is $u'\mapsto (u')^b$.  On the free boundary
$q=0$, the forward map is $u\mapsto u^b$ and the reverse map is
$u'\mapsto (u')^{1/b}$.

For terminal parameter $t$, define the free density
\begin{equation}\label{eq:free-a}
 a_h(t):=1-e^{-t/b^h}.
\end{equation}
The exact free orbit with this density is
\begin{equation}\label{eq:free-orbit}
 \bar u_i=e^{-t/b^{h-i}},\qquad
 \bar\rho_i=1-\bar u_i,\qquad
 \bar q_i=0.
\end{equation}


\section{Uniform one-step estimates}
The slack theorem requires only terminal parameters linear in $h$ and uniformly below the critical coefficient.

\begin{definition}[Subcritical terminal window]\label{def:terminal-window}
Fix constants
\[
 0<\tau_-\le\tau_+<\frac{4\log b}{3D}.
\]
A sequence of terminal parameters is in the subcritical linear window if
\[
 \tau_-h\le t\le\tau_+h.
\]
\end{definition}

All constants below may depend on $d,\tau_-,\tau_+$ but not on $h$ or $t$ in this window.

\subsection{Free-side perturbations}
At one reverse step, keep $\rho'$ fixed and write $q'=1-v'$.  The free predecessor with the same $\rho'$ has
\[
 w=(1-\rho')^{1/b},
 \qquad
 \rho_0=1-w,
 \qquad
 u_0=w,
 \qquad
 q_0=0.
\]

\begin{lemma}[Uniform free-side step]\label{lem:free-step}
There are $\delta_0,C>0$ such that, whenever $0\le q'\le\delta_0$, the actual predecessor satisfies
\begin{align}
 1-\rho&=w(1+\eta_\rho),\label{eq:free-rho-step}\\
 q&=bw^2q'(1+\eta_q),\label{eq:free-q-step}
\end{align}
with
\[
 |\eta_\rho|+|\eta_q|\le Cq'.
\]
The estimates are uniform for $0<w\le1$.
\end{lemma}

\begin{proof}
Put $\delta=q'$.  Since $v'=1-\delta$,
\[
 R=(1-\delta)\frac{1-w}{w},
 \qquad
 M=\left(1-w+\frac w{1-\delta}\right)^b
 =\left(1+\frac{w\delta}{1-\delta}\right)^b.
\]
For $\delta\le1/2$, the binomial expansion with a uniform remainder gives
\[
 M=1+bw\delta+O_d(\delta^2).
\]
Furthermore,
\[
 w(R+M)
 =(1-\delta)(1-w)+wM
 =1+O_d(\delta).
\]
Now
\[
 1-\rho=u+q=\frac M{R+M},
 \qquad
 q=\frac{M-1}{R+M}.
\]
Substitution gives~\eqref{eq:free-rho-step}--\eqref{eq:free-q-step}.
\end{proof}

\subsection{Stable-side perturbations}
For a reverse step near the stable manifold, write the next state as
\[
 v'=V,
 \qquad
 \rho'=V\delta,
 \qquad
 u'=V(1-\delta).
\]
The stable predecessor with the same $V$ is $(\rho_0,u_0)=(0,V^b)$.

\begin{lemma}[Uniform stable-side step]\label{lem:stable-step}
There are $\delta_1,C>0$ such that, whenever $0\le\delta\le\delta_1$, the predecessor satisfies
\begin{align}
 \log u&=b\log V+O_d(V\delta),\label{eq:stable-logu}\\
 \frac\rho v&=\frac{V^2}{b}\delta(1+O_d(\delta)).\label{eq:stable-delta}
\end{align}
Put
\[
 \widetilde a_i:=\frac{\rho_i}{b^iu_i^D}.
\]
Then
\begin{equation}\label{eq:stable-invariant-step}
 \frac{\widetilde a_i}{\widetilde a_{i+1}}
 =1+O_d(\delta).
\end{equation}
For sufficiently small $\delta_1$, the transverse ratios contract backward:
\[
 \frac{\rho_i}{v_i}\le\frac{1+o(1)}b\frac{\rho_{i+1}}{v_{i+1}}.
\]
\end{lemma}

\begin{proof}
Here
\[
 w=(1-V\delta)^{1/b},
 \qquad
 e=1-w=\frac{V\delta}{b}(1+O_d(\delta)),
\]
so
\[
 R=V\frac ew=\frac{V^2\delta}{b}(1+O_d(\delta)).
\]
Also
\[
 e+\frac wV
 =\frac1V\left(w+Ve\right)
 =\frac1V\left(1-e(1-V)\right),
\]
whence
\[
 M=V^{-b}(1+O_d(V\delta)).
\]
Since $R/M=O_d(V^{b+2}\delta)$, equation~\eqref{eq:reverse-out} yields~\eqref{eq:stable-logu}; moreover $\rho/u=R$, which gives~\eqref{eq:stable-delta}.

For~\eqref{eq:stable-invariant-step}, use $\rho=uR$ and compute
\[
 \frac{\widetilde a_i}{\widetilde a_{i+1}}
 =\frac{bRu^{1-D}(u')^D}{\rho'}.
\]
The preceding expansions give
\[
 \frac{bR}{\rho'}=V(1+O_d(\delta)),
 \qquad
 u^{1-D}(u')^D=V^{-1}(1+O_d(\delta)),
\]
because $(b-1)D=b+1$.  Their product is $1+O_d(\delta)$.
\end{proof}


\section{Free shadowing to an overlap layer}
In the quantitative statements below, notation of the form
$1+B_h^{-\eta+o(1)}$ denotes $1+\epsilon_h$ with
$|\epsilon_h|\le B_h^{-\eta+o(1)}$; no sign is implied.
For $t$ in the terminal window, choose
\begin{equation}\label{eq:H-choice}
 H=\left\lfloor\frac{3Dt}{4\log b}\right\rfloor,
 \qquad
 m=h-H.
\end{equation}
Then $H,m$ are both linear in $h$.  Define the formal free transverse sequence, for $m\le i\le h$, by
\begin{equation}\label{eq:qbar}
 \bar q_i
 =b^{h-i}
 \exp\left\{-Dt+\frac{2t}{b-1}b^{i-h}\right\}.
\end{equation}
It satisfies $\bar q_h=e^{-t}$ and
\[
 \bar q_i=b\bar u_i^2\bar q_{i+1}.
\]

\begin{proposition}[Uniform free shadow]\label{prop:free-shadow}
Uniformly in every subcritical terminal window, for every $m\le i\le h$,
\begin{align}
 q_i&=\bar q_i\left(1+O\!\left(h e^{-Dt/4}\right)\right),\label{eq:q-shadow}\\
 1-\rho_i&=\bar u_i\left(1+O\!\left(e^{-Dt/4}\right)\right).\label{eq:rho-shadow}
\end{align}
In particular,
\begin{align}
 q_m&=b^H\exp\left\{-Dt+\frac{2t}{b-1}b^{-H}\right\}
 \left(1+O\!\left(h e^{-Dt/4}\right)\right),\label{eq:qm}\\
 \rho_m&=t b^{-H}
 \left(1+O\!\left(t b^{-H}+e^{-Dt/4}\right)\right),\label{eq:rhom}\\
 \frac{\rho_m}{q_m}&=\exp\{-Dt/2+o(h)\}.\label{eq:rho-over-q}
\end{align}
If $L=\log B_h$ and $c_t=Dt/L$, then, uniformly when $c_t$ stays in a compact subset of the terminal window,
\begin{equation}\label{eq:overlap-scales-quantitative}
 q_m=B_h^{-c_t/4+o(1)},\qquad
 \rho_m=B_h^{-3c_t/4+o(1)},\qquad
 \frac{\rho_m}{q_m}=B_h^{-c_t/2+o(1)}.
\end{equation}
\end{proposition}

\begin{proof}
Let
\[
 E_i=\log\frac{1-\rho_i}{\bar u_i}.
\]
Lemma~\ref{lem:free-step} gives
\[
 E_i=\frac1bE_{i+1}+O(q_{i+1}),
\]
and
\[
 \log\frac{q_i}{\bar q_i}
 =\log\frac{q_{i+1}}{\bar q_{i+1}}
 +\frac2bE_{i+1}+O(q_{i+1}).
\]
Both errors vanish at $i=h$.

Write $r=h-i$.  The logarithm of the formal transverse sequence is
\[
 f(r)=r\log b-Dt+\frac{2t}{b-1}b^{-r}.
\]
The function $f$ is convex, so its maximum on $0\le r\le H$ occurs at an endpoint.  At $r=0$, $f(0)=-t$, while the definition of $H$ gives
\[
 f(H)\le-\frac14Dt+O_d(1).
\]
Since $D/4<1$, it follows that
\[
 \max_{m\le i\le h}\bar q_i=O_d(e^{-Dt/4}).
\]

A first-failure bootstrap now gives
\[
 \max_i|E_i|=O(e^{-Dt/4}),
 \qquad
 \max_i\left|\log\frac{q_i}{\bar q_i}\right|
 =O(h e^{-Dt/4}).
\]
For large $h$ these estimates prevent the first failure, proving
\eqref{eq:q-shadow}--\eqref{eq:rho-shadow} and the stated formula for $q_m$.
Furthermore,
\[
 1-e^{-t/b^H}=t b^{-H}\left(1+O(t b^{-H})\right),
\]
which gives~\eqref{eq:rhom}.  Finally,
\[
 \log\frac{\rho_m}{q_m}
 =\log t-2H\log b+Dt+o(h)
 =-\frac12Dt+o(h).
\]
Because $H=3Dt/(4\log b)+O(1)$ and $L=h\log b+O(1)$, the three estimates in~\eqref{eq:overlap-scales-quantitative} follow.
\end{proof}


\section{Stable shadowing from the overlap to the root}
Put
\[
 x=-\log u_1,
 \qquad
 \widetilde a=\widetilde a_1
 =\frac{\rho_1}{b u_1^D}.
\]

\begin{proposition}[Uniform stable shadow]\label{prop:stable-shadow}
Under the terminal-window assumptions and the matching choice~\eqref{eq:H-choice}, put
\[
 \chi_m:=q_m+\frac{\rho_m}{q_m}.
\]
Then
\begin{align}
 x&=b^{m-1}q_m\left(1+O(\chi_m)\right),\label{eq:x-from-q}\\
 \widetilde a&=\frac{\rho_m}{b^m}
 \left(1+O(q_m+\rho_m)\right).\label{eq:atilde-from-rho}
\end{align}
If the messages are normalized by $A_1=1$, then
\begin{equation}\label{eq:Am-stable}
 \log A_m
 =\frac{b}{b-1}\left(1-b^{-(m-1)}\right)\log u_1
 +O(m\rho_m),
\end{equation}
and
\begin{equation}\label{eq:r-Am}
 B_0=\widetilde a\,A_m^2
 \left(1+O(q_m+m\rho_m)\right).
\end{equation}
All constants are uniform in the terminal window.
\end{proposition}

\begin{proof}
Set $\delta_i=\rho_i/v_i$.  Proposition~\ref{prop:free-shadow} gives $\delta_m=O(\rho_m)$, and Lemma~\ref{lem:stable-step} implies geometric backward contraction.  Hence
\[
 \sum_{i=2}^m\delta_i=O(\rho_m).
\]
Iteration of~\eqref{eq:stable-invariant-step} gives
\[
 \widetilde a
 =\frac{\rho_m}{b^m u_m^D}\left(1+O(\rho_m)\right).
\]
Since $u_m=1-q_m-\rho_m$, this proves~\eqref{eq:atilde-from-rho}.

The stable-step estimate also gives
\[
 \log v_i=b\log v_{i+1}+O(\rho_{i+1}).
\]
After iteration,
\[
 -\log v_1
 =b^{m-1}[-\log v_m]
 +O\left(\sum_{j=2}^m b^{j-2}\rho_j\right)
 =b^{m-1}[-\log v_m]+O(b^{m-1}\rho_m).
\]
Since $v_m=1-q_m$,
\[
 -\log v_m=q_m(1+O(q_m)).
\]
Backward contraction gives $\delta_1=O(b^{-(m-1)}\rho_m)$, and therefore
\[
 0\le\log\frac{v_1}{u_1}
 =\log\left(1+\frac{\rho_1}{u_1}\right)
 =O(b^{-(m-1)}\rho_m).
\]
This is absorbed by the preceding error.  Dividing by the main term
$b^{m-1}q_m$ proves~\eqref{eq:x-from-q}.

Likewise,~\eqref{eq:stable-logu} and contraction give
\[
 \log u_i=b^{-(i-1)}\log u_1+O(\rho_m)
\]
uniformly for $1\le i<m$.  Summing the message ratios gives~\eqref{eq:Am-stable}.

At the root,
\[
 B_0=v_1\frac{1-(1-\rho_1)^{1/b}}{(1-\rho_1)^{1/b}}
 =\widetilde a\,u_1^{D+1}\left(1+O(\rho_m)\right).
\]
Since $D+1=2b/(b-1)$, equations~\eqref{eq:x-from-q} and~\eqref{eq:Am-stable} give
\[
 \log\frac{B_0}{\widetilde a A_m^2}
 =-\frac{2b}{b-1}\frac{x}{b^{m-1}}+O(m\rho_m)
 =O(q_m+m\rho_m),
\]
which proves~\eqref{eq:r-Am}.
\end{proof}


\section{Density and activity asymptotics}
The free and stable shadows now match without an additional hypothesis.

\begin{theorem}[Terminal-window matching]\label{thm:matching}
Let $L=\log B_h$ and $c_t=Dt/L$.  Uniformly for $t$ in every subcritical linear window,
\begin{align}
 \widetilde a&=a_h(t)\left(1+B_h^{-c_t/4+o(1)}\right),\label{eq:atilde-a}\\
 \frac{x}{b^{h-1}}
 &=(1-a_h(t))^{B_h}
 \left(1+B_h^{-c_t/4+o(1)}\right),\label{eq:x-coupon}\\
 B_h\lvert\widetilde a-a_h(t)\rvert
 &=B_h^{-c_t/4+o(1)}.\label{eq:atilde-a-absolute}
\end{align}
\end{theorem}

\begin{proof}
Equations~\eqref{eq:rhom}, \eqref{eq:atilde-from-rho}, and
\eqref{eq:overlap-scales-quantitative} give
\[
 \widetilde a=\frac{t}{b^h}
 \left(1+B_h^{-c_t/4+o(1)}\right)
 =a_h(t)\left(1+B_h^{-c_t/4+o(1)}\right).
\]
Since $a_h(t)B_h=Dt(1+o(1))=B_h^{o(1)}$, this also proves~\eqref{eq:atilde-a-absolute}.

Combining~\eqref{eq:qm} and~\eqref{eq:x-from-q} yields
\[
 \frac{x}{b^{h-1}}
 =\exp\left\{-Dt+\frac{2t}{b-1}b^{-H}\right\}
 \left(1+B_h^{-c_t/4+o(1)}\right).
\]
On the other hand, by~\eqref{eq:Bh} and~\eqref{eq:free-a},
\[
 (1-a_h(t))^{B_h}
 =\exp\left\{-Dt+\frac{2t}{b-1}b^{-h}\right\}.
\]
The exponent difference is $O(tb^{-H})=B_h^{-3c_t/4+o(1)}$, which is smaller than the displayed error.  This proves~\eqref{eq:x-coupon}.
\end{proof}

To pass from $\widetilde a$ to the actual selected density, reconstruct the edge partition function
\[
 Z_e=\sum_{i=0}^hB_i^2+2\sum_{i=0}^{h-1}B_iA_{i+1}.
\]

\begin{lemma}[Edge-normalizer localization]\label{lem:Ze-localization}
With the matching index $m$,
\begin{equation}\label{eq:Ze-Am}
 Z_e=A_m^2\left(1+O(hq_m+\rho_m)\right).
\end{equation}
Consequently, with $c_t=Dt/\log B_h$,
\begin{align}
 \alpha&=\widetilde a\left(1+B_h^{-c_t/4+o(1)}\right)
 =a_h(t)\left(1+B_h^{-c_t/4+o(1)}\right),\label{eq:alpha-atilde}\\
 B_h|\alpha-a_h(t)|&=B_h^{-c_t/4+o(1)}.\label{eq:alpha-free-absolute}
\end{align}
\end{lemma}

\begin{proof}
For $i\ge m$,
\[
 B_i^2+2B_iA_{i+1}
 =(B_i+A_{i+1})^2-A_{i+1}^2.
\]
Since $B_i+A_{i+1}=A_i(1-q_i)$, summation gives the exact tail identity
\[
 \sum_{i=m}^h(B_i^2+2B_iA_{i+1})
 =A_m^2(1-q_m)^2
 -\sum_{i=m+1}^hA_i^2q_i(2-q_i).
\]
Proposition~\ref{prop:free-shadow} bounds the relative tail error by $O(hq_m)$.

For the left part, Proposition~\ref{prop:stable-shadow} and the stable invariant give, uniformly for $i<m$,
\[
 \rho_i=\widetilde a\,b^iu_i^D(1+O(\rho_m)),
\]
while
\[
 \frac{A_i^2u_i^{D+1}}{A_m^2}=1+O(q_m+m\rho_m).
\]
Therefore
\[
 \frac1{A_m^2}
 \sum_{i=1}^{m-1}2B_iA_{i+1}
 =O\!\left(\widetilde a\sum_{i=1}^{m-1}b^i\right)
 =O(\widetilde a b^m)
 =O(\rho_m).
\]
The square terms are smaller by a factor $\rho_i/u_i$, and the root term is $O(\widetilde a A_m^2)$.  This proves~\eqref{eq:Ze-Am}.

The exact root belief is
\[
 \alpha=\frac{B_0(B_0+A_1)}{Z_e}.
\]
Use $A_1=1$,~\eqref{eq:r-Am},~\eqref{eq:Ze-Am}, and $B_0=o(1)$ to obtain
\[
 \alpha=\widetilde a\left(1+O(hq_m+\rho_m)\right).
\]
Now apply Theorem~\ref{thm:matching} and
\eqref{eq:overlap-scales-quantitative}.  Polynomial factors in $h$ are
$B_h^{o(1)}$, which proves~\eqref{eq:alpha-atilde} and
\eqref{eq:alpha-free-absolute}.
\end{proof}

The root reconstruction formulas are
\[
 B_0=v_1\frac{1-(1-\rho_1)^{1/b}}{(1-\rho_1)^{1/b}},
 \qquad
 \kappa=\left(\frac{v_1}{(1-\rho_1)^{1/b}}\right)^b,
 \qquad
 z=\frac{B_0\kappa}{(1+B_0)^b}.
\]

\begin{theorem}[Quantitative uniform activity law]\label{thm:activity-law}
Fix constants $0<c_-<c_+<4/3$.  Uniformly for densities satisfying
\begin{equation}\label{eq:density-window}
 c_-\log B_h\le\alpha B_h\le c_+\log B_h,
 \qquad
 c=\frac{\alpha B_h}{\log B_h},
\end{equation}
one has
\begin{equation}\label{eq:activity-additive-quantitative}
 -\log z(\alpha)-\log\frac1\alpha
 =B_h(1-\alpha)^{B_h}
  \left(1+B_h^{-c/4+o(1)}\right)
  +B_h^{-c/4+o(1)}.
\end{equation}
Consequently,
\begin{equation}\label{eq:activity-law}
 \boxed{
 -\log z(\alpha)
 =\log\frac1\alpha
 +B_h(1-\alpha)^{B_h}
 \left(1+B_h^{-\delta+o(1)}\right),
 }
\end{equation}
where
\begin{equation}\label{eq:activity-delta}
 \delta=\min\left\{\frac{c_-}{4},
 1-\frac{3c_+}{4}\right\}>0.
\end{equation}
In particular, for all sufficiently large $h$,
\begin{equation}\label{eq:activity-lower}
 -\log z(\alpha)
 \ge\log\frac1\alpha
 +\frac12B_h(1-\alpha)^{B_h}
\end{equation}
uniformly on~\eqref{eq:density-window}.
\end{theorem}

\begin{proof}
First work in a terminal window and put $c_t=Dt/\log B_h$.  We spell out the quantitative root reconstruction because its additive error is later compared with a vanishing coupon term.  Proposition~\ref{prop:stable-shadow} and the exact root formulas give
\[
 B_0=\widetilde a\,u_1^{D+1}
 \left(1+O(q_m+m\rho_m)\right),
 \qquad
 \kappa=u_1^b\left(1+O(\rho_m)\right).
\]
Since $z=B_0\kappa/(1+B_0)^b$, $x=-\log u_1$, and
$B_0=B_h^{-1+o(1)}$, it follows that
\[
 -\log z
 =\log\frac1{\widetilde a}+Kx
 +O(q_m+m\rho_m+B_0)
 =\log\frac1{\widetilde a}+Kx
 +B_h^{-c_t/4+o(1)},
 \qquad K=b+D+1=bD.
\]
The arithmetic identity
\[
 Kb^{h-1}=Db^h=B_h+\frac2{b-1}
\]
and Theorem~\ref{thm:matching} yield
\[
 Kx
 =B_h(1-a_h(t))^{B_h}
 \left(1+B_h^{-c_t/4+o(1)}\right).
\]
Lemma~\ref{lem:Ze-localization} gives
\[
 \log\frac{\alpha}{\widetilde a}=B_h^{-c_t/4+o(1)},
 \qquad
 B_h|\alpha-a_h(t)|=B_h^{-c_t/4+o(1)}.
\]
The latter estimate changes the coupon term relatively by
$1+B_h^{-c_t/4+o(1)}$.  Hence
\begin{equation}\label{eq:activity-terminal-quantitative}
 -\log z-\log\frac1\alpha
 =B_h(1-\alpha)^{B_h}
  \left(1+B_h^{-c_t/4+o(1)}\right)
  +B_h^{-c_t/4+o(1)}.
\end{equation}

It remains to identify the terminal range corresponding to~\eqref{eq:density-window}.  By~\eqref{eq:alpha-atilde}, uniformly in every terminal window,
\[
 \alpha B_h=Dt\left(1+B_h^{-c_t/4+o(1)}\right).
\]
As in the qualitative proof, choose a slightly broader terminal window whose endpoint densities bracket~\eqref{eq:density-window}; strict activity and density monotonicity then expose every density in the displayed interval.  Uniformly there, $c_t=c+o(1)$, and~\eqref{eq:activity-terminal-quantitative} becomes~\eqref{eq:activity-additive-quantitative}.

Finally,
\[
 B_h(1-\alpha)^{B_h}=B_h^{1-c+o(1)}.
\]
The additive error in~\eqref{eq:activity-additive-quantitative} is therefore relatively
$B_h^{3c/4-1+o(1)}$, while the multiplicative orbit error is
$B_h^{-c/4+o(1)}$.  Taking the worst exponent over $[c_-,c_+]$ proves
\eqref{eq:activity-law}--\eqref{eq:activity-delta};
\eqref{eq:activity-lower} follows immediately.
\end{proof}


\section{A capped-free profile and the bounded critical window}\label{sec:bounded-window}
Put
\[
 B=B_h,\qquad L=\log B,
 \qquad
 \Dcal_h(\alpha):=H(\alpha)-\Psi_{d,h}(\alpha),
 \qquad
 Q_h(\alpha):=(1-\alpha)^B.
\]
We now construct a feasible profile whose only entropy loss is the conditioning event at the terminal wall.

Put $r=1-\alpha$ and
\[
 T_j=1+b+\cdots+b^j=\frac{b^{j+1}-1}{b-1}.
\]
For $h\ge2$, define
\begin{equation}\label{eq:capped-free-messages}
 A_i=r^{T_{i-1}}\quad(1\le i\le h),
 \qquad A_{h+1}=0,
\end{equation}
\begin{equation}\label{eq:capped-free-B}
 B_0=\alpha,
 \qquad B_i=A_i-A_{i+1}\quad(1\le i<h),
 \qquad B_h=A_h,
\end{equation}
and set
\begin{equation}\label{eq:capped-free-beliefs}
 \ell_i=B_i^2,
 \qquad
 x_i=B_{i-1}A_i.
\end{equation}

\begin{lemma}[Capped-free normalization and feasibility]\label{lem:capped-free-feasible}
The profile~\eqref{eq:capped-free-beliefs} is normalized and has selected density $p_0=\alpha$.  For $i<h$, its lower-neighbor marginal is the nonempty conditioned-binomial marginal with tilt $B_{i-1}/A_i$ and is strictly larger than $1/d$.  At the terminal level,
\begin{equation}\label{eq:capped-terminal}
 p_h=A_{h-1}A_h,
 \qquad
 a_h=1-\varepsilon_h,
 \qquad
 \varepsilon_h:=\frac{A_h}{A_{h-1}}=r^{b^{h-1}}.
\end{equation}
Consequently the profile is feasible whenever
$\varepsilon_h\le(d-1)/d$.

More explicitly, for the window $|s|\le M$ in
Theorem~\ref{thm:bounded-critical-window}, feasibility holds for every
$h\ge h_0(d,M)$, where
\begin{equation}\label{eq:explicit-feasibility-threshold}
 h_0(d,M):=\min\left\{h\ge2:
 L_h-2\log L_h-M
 \ge \frac{d(d-1)}{d-2}\log\frac d{d-1}
 \right\}.
\end{equation}
\end{lemma}

\begin{proof}
Since $B_0+A_1=1$ and $B_i+A_{i+1}=A_i$ for $i\ge1$,
\[
 B_0^2+2B_0A_1=1-A_1^2,
 \qquad
 B_i^2+2B_iA_{i+1}=A_i^2-A_{i+1}^2.
\]
Summing gives total directed-edge mass one, and
$p_0=B_0(B_0+A_1)=\alpha$.

For $i<h$, put $S_i=B_{i-1}+B_i+A_{i+1}$.  The messages satisfy
\[
 \kappa A_i=S_i^b,
 \qquad
 \kappa B_i=S_i^b-A_i^b,
 \qquad
 \kappa=1/r.
\]
It follows that
\[
 \kappa p_i=S_i^d-A_i^d,
 \qquad
 \kappa x_i=B_{i-1}S_i^b,
\]
which is exactly the conditioned-binomial row identity with tilt
$B_{i-1}/A_i$.  A positive tilt gives expected nonempty-subset size strictly larger than one, hence lower-neighbor marginal strictly larger than $1/d$.  The terminal formulas follow directly from
$B_h=A_h$ and $B_{h-1}=A_{h-1}-A_h$.

It remains to verify the explicit threshold.  Since
$B_h<Db^h$,
\[
 \alpha b^{h-1}
 >\frac{C}{Db}
 =\frac{d-2}{d(d-1)}C.
\]
Under~\eqref{eq:explicit-feasibility-threshold} and $|s|\le M$, this is at least $\log(d/(d-1))$.  Therefore
\[
 \varepsilon_h=(1-\alpha)^{b^{h-1}}
 \le e^{-\alpha b^{h-1}}
 \le\frac{d-1}{d},
\]
as required.
\end{proof}

Define
\[
 \Delta_d(\varepsilon)
 :=dH(\varepsilon)-s_d(1-\varepsilon).
\]

\begin{lemma}[Exact capped-free entropy loss]\label{lem:capped-free-loss}
Let $\Fcal_{d,h}^{\rm cap}(\alpha)$ be the compact functional at the profile
\eqref{eq:capped-free-beliefs}.  Then
\begin{equation}\label{eq:capped-free-loss}
 H(\alpha)-\Fcal_{d,h}^{\rm cap}(\alpha)
 =p_h\Delta_d(\varepsilon_h),
\end{equation}
and
\begin{equation}\label{eq:capped-coupon-exact}
 p_h\varepsilon_h^d=(1-\alpha)^B.
\end{equation}
Moreover, as $\varepsilon\downarrow0$,
\begin{equation}\label{eq:terminal-gap-asymptotic}
 \Delta_d(\varepsilon)
 =\varepsilon^d\left(1+O_d(\varepsilon^{d-1})\right).
\end{equation}
\end{lemma}

\begin{proof}
First replace the terminal entropy $s_d(a_h)$ by the unconstrained coordinate entropy $dH(a_h)=dH(\varepsilon_h)$.  For every nonterminal layer, and for the relaxed terminal layer, the dual entropy identity and the message row identities reduce the layer contribution to
\[
 p_i\log\kappa
 -d x_i\log B_{i-1}
 +d x_{i+1}\log B_i,
 \qquad x_{h+1}=0.
\]
The sum over $i=1,\ldots,h$ telescopes to
\[
 (1-\alpha)\log\kappa-dx_1\log B_0.
\]
The root contribution is
\[
 (d-1)\alpha\log\alpha
 -\frac d2\alpha^2\log(\alpha^2).
\]
Using $\kappa=1/r$, $x_1=\alpha r$, and $B_0=\alpha$, the relaxed total is exactly $H(\alpha)$.  Restoring the nonempty terminal condition subtracts the right side of~\eqref{eq:capped-free-loss}.

Since $A_h=rA_{h-1}^b$,
\[
 p_h\varepsilon_h^d
 =A_{h-1}A_h\left(\frac{A_h}{A_{h-1}}\right)^d
 =rA_h^d
 =r^{1+dT_{h-1}}
 =r^B,
\]
which proves~\eqref{eq:capped-coupon-exact}.

For~\eqref{eq:terminal-gap-asymptotic}, let $\eta$ be the coordinate-exclusion probability before conditioning on a nonempty subset.  Then
\[
 \varepsilon=\frac{\eta-\eta^d}{1-\eta^d},
 \qquad
 1-\varepsilon=\frac{1-\eta}{1-\eta^d},
 \qquad
 \eta-\varepsilon=(1-\varepsilon)\eta^d,
\]
so $\eta=\varepsilon+O(\varepsilon^d)$.  Substituting the first two identities into the entropy of the conditioned product measure gives
\[
 \Delta_d(\varepsilon)
 =-\log(1-\eta^d)
 -d\,\mathrm{KL}\!\left(
 \operatorname{Ber}(\varepsilon)
 \Vert\operatorname{Ber}(\eta)
 \right).
\]
The Bernoulli chi-square bound makes the KL term
$O_d(\varepsilon^{2d-1})$, while
$-\log(1-\eta^d)=\varepsilon^d(1+O_d(\varepsilon^{d-1}))$.
\end{proof}

\begin{proof}[Proof of Theorem~\ref{thm:bounded-critical-window}]
The upper type anchor gives $\Psi_{d,h}(\alpha)\le H(\alpha)$, while
Lemma~\ref{lem:capped-free-feasible} supplies a feasible compact profile for every sufficiently large $h$, uniformly in the displayed window.  Hence Lemma~\ref{lem:capped-free-loss} implies
\[
 0\le \Dcal_h(\alpha)
 \le Q_h(\alpha)
 \left(1+O_d(\varepsilon_h^{d-1})\right).
\]
There,
\[
 \varepsilon_h=(1-\alpha)^{b^{h-1}}
 =B^{-\frac{d-2}{d(d-1)}+o(1)},
\]
which proves the upper bound in~\eqref{eq:defect-bounded-two-sided}.

For the lower bound, define
\[
 \mathcal A_h(a)
 :=-\log z(a)-\log\frac1a.
\]
The envelope identity~\eqref{eq:globality} gives
\begin{equation}\label{eq:defect-derivative}
 -\Dcal_h'(a)=-\log(1-a)+\mathcal A_h(a).
\end{equation}
Put
\[
 C_0=L-2\log L+s,
 \quad \alpha_0=C_0/B,
 \qquad
 C_1=\frac87L,
 \quad \alpha_1=C_1/B.
\]
For all large $h$, the interval lies in the fixed density window with
$c_-=3/4$ and $c_+=8/7$.  Theorem~\ref{thm:activity-law} applies uniformly there with relative error $B^{-1/7+o(1)}$.  Since
$-\log(1-a)\ge0$ and $\Dcal_h(\alpha_1)\ge0$,
\begin{align*}
 \Dcal_h(\alpha_0)
 &\ge
 \left(1-B^{-1/7+o(1)}\right)
 B\int_{\alpha_0}^{\alpha_1}(1-a)^B\,da\\
 &=\left(1-B^{-1/7+o(1)}\right)
 \frac{B}{B+1}
 \left[(1-\alpha_0)^{B+1}-(1-\alpha_1)^{B+1}\right].
\end{align*}
Finally,
\[
 \frac{(1-\alpha_1)^B}{(1-\alpha_0)^B}
 =B^{-1/7}L^{-2}e^{s+o(1)},
\]
and $\alpha_0=O(L/B)$.  This proves the lower bound in
\eqref{eq:defect-bounded-two-sided} and hence~\eqref{eq:defect-bounded-window}.
\end{proof}

\begin{proof}[Proof of Corollary~\ref{cor:critical-scaling-function}]
Uniformly for bounded $s$,
\[
 H\!\left(\frac{L-2\log L+s}{B}\right)
 =\frac{L^2}{B}(1+o(1)),
\]
whereas
\[
 \left(1-\frac{L-2\log L+s}{B}\right)^B
 =e^{-s}\frac{L^2}{B}(1+o(1)).
\]
Apply Theorem~\ref{thm:bounded-critical-window}.
\end{proof}

\begin{proof}[Proof of Corollary~\ref{cor:annealed-zero-expansion}]
For every fixed $\epsilon>0$, Corollary~\ref{cor:critical-scaling-function} gives opposite signs at $s=-\epsilon$ and $s=\epsilon$ for all large $h$; because strict concavity makes the superlevel set
$\{\alpha:\Psi_{d,h}(\alpha)\ge0\}$ an interval, its lower endpoint satisfies
\[
 C_h^{\rm ann}=L-2\log L+o(1).
\]

Let
\[
 F_B(C):=H(C/B)-(1-C/B)^B.
\]
Theorem~\ref{thm:bounded-critical-window} gives
\[
 \Psi_{d,h}(C/B)
 =F_B(C)+O\!\left((1-C/B)^B B^{-1/7+o(1)}\right)
\]
uniformly in a fixed neighborhood of the scalar root.  There,
\[
 F_B'(C)=(1+o(1))(1-C/B)^B,
\]
uniformly and with positive sign.  A two-sided mean-value comparison therefore proves~\eqref{eq:graph-scalar-zero}.

It remains to invert the scalar equation.  For $C=L+O(\ell)$,
\eqref{eq:scalar-coupon-root} is equivalent to
\begin{equation}\label{eq:scalar-log-equation}
 C+\log C+\log(L-\log C+1)
 =L+O(L^2/B).
\end{equation}
The equation first gives $C=L-2\ell+O(\ell/L)$; substituting this estimate successively into the two logarithms determines the coefficients through order $L^{-3}$.  If $P_3$ denotes the displayed polynomial in~\eqref{eq:annealed-zero-expansion}, direct expansion gives
\[
 P_3+\log P_3+\log(L-\log P_3+1)-L
 =O(\ell^4/L^4),
\]
while the derivative of the left side of~\eqref{eq:scalar-log-equation} is $1+O(1/L)$.  Thus the mean-value theorem yields
\[
 \widehat C_B-P_3
 =O(\ell^4/L^4+L^2/B)
 =O(\ell^4/L^4),
\]
which supplies the claimed bootstrap remainder.  Finally,
$B^{-1/7+o(1)}$ is smaller than every fixed inverse power of $L$, so~\eqref{eq:graph-scalar-zero} transfers the expansion to $C_h^{\rm ann}$.  Repeating the same formal substitution and mean-value estimate yields the expansion to any prescribed fixed inverse power of $L$.
\end{proof}

\section{Near-critical integration and the random-graph deduction}\label{sec:integration}
We first prove Theorem~\ref{thm:near-critical}.  The pointwise anchor is Corollary~\ref{cor:Psi-le-H}.

\begin{proof}[Proof of Theorem~\ref{thm:near-critical}]
Put $B=B_h$ and $L=\log B$.  For a value $C$ in~\eqref{eq:C-window}, set
\[
 \alpha_- =\frac CB,
 \qquad
 C^\sharp=\frac{L+C}{2},
 \qquad
 \alpha_+=\frac{C^\sharp}{B}.
\]
For large $h$, the interval $[\alpha_-,\alpha_+]$ lies in a fixed density window
\[
 \frac\eta2\frac LB
 \le a\le
 \frac LB,
\]
so Theorem~\ref{thm:activity-law} applies uniformly.  Exact globality and the envelope identity~\eqref{eq:globality} give
\[
 \Psi(\alpha_-)
 =\Psi(\alpha_+)+\int_{\alpha_-}^{\alpha_+}\log z(a)\,da.
\]
Corollary~\ref{cor:Psi-le-H} and~\eqref{eq:activity-lower} imply
\begin{align*}
 \Psi(\alpha_-)
 &\le H(\alpha_+)
 -\int_{\alpha_-}^{\alpha_+}\log\frac1a\,da
 -\frac12\int_{\alpha_-}^{\alpha_+}B(1-a)^B\,da\\
 &\le H(\alpha_-)
 -\frac12\frac{B}{B+1}
 \left[(1-\alpha_-)^{B+1}-(1-\alpha_+)^{B+1}\right].
\end{align*}
The second inequality uses $H'(a)=\log(1/a)+\log(1-a)$.

Write
\[
 W(C):=L-2\log L-C.
\]
By hypothesis $W(C)\ge W_h\to\infty$.  Uniformly over~\eqref{eq:C-window},
\[
 H(\alpha_-)=O\left(\frac{L^2}{B}\right)
 =o(e^{-C}),
\]
because
\[
 e^{-C}=e^{W(C)}\frac{L^2}{B}.
\]
Also $C=O(L)$ gives
\[
 (1-\alpha_-)^{B+1}=e^{-C}(1+o(1)).
\]
Finally,
\[
 C^\sharp-C=\frac{L-C}{2}
 \ge\log L+\frac{W_h}{2},
\]
so
\[
 (1-\alpha_+)^{B+1}
 =e^{-C^\sharp}(1+o(1))
 =o(e^{-C}).
\]
All estimates are uniform, proving~\eqref{eq:near-critical-psi}.
\end{proof}

\begin{proof}[Proof of Corollary~\ref{cor:fixed-slack}]
For $C=cL$ with $c\in[c_-,c_+]$, one has
\[
 L-2\log L-C=(1-c)L-2\log L\to\infty
\]
uniformly.  Theorem~\ref{thm:near-critical} gives $e^{-C}=B^{-c}$.
\end{proof}

\begin{proof}[Proof of Theorem~\ref{thm:random-fixed-window}]
Put
\[
 C_h^*=L_h-2\log L_h-\omega,
 \qquad
 T_n=\frac{nC_h^*}{B_h}.
\]
If $T_n\le1$, then the real-valued lower bound follows from
$\gamma_h(G)\ge1$.  Otherwise set
\[
 m_n=\lceil T_n\rceil-1,
 \qquad
 C_n=\frac{m_nB_h}{n}.
\]
Condition~\eqref{eq:fixed-window-growth} implies $B_h/n=o(1)$.  Indeed, $h=\Theta_d(L_h)$.  If $B_h\ge n$ along a subsequence and $x=B_h/n\ge1$, then
\[
 \frac{nL_h^2/B_h}{h\log n}
 =\Theta_d\!\left(
 \frac{L_h}{x\log n}
 \right)
 =\Theta_d\!\left(
 \frac{1+\log x/\log n}{x}
 \right)
 =O_d(1),
\]
contradicting~\eqref{eq:fixed-window-growth}.  Hence $B_h<n$ eventually, so $L_h\le\log n$; the same displayed ratio tends to infinity only if $n/B_h\to\infty$.  Therefore
\[
 C_n=C_h^*+o(1),
 \qquad C_n\le C_h^*.
\]
The scaling-function convergence is uniform on compact $s$-intervals; hence
\[
 \Psi_{d,h}(C_n/B_h)
 =-\left(e^\omega-1+o(1)\right)
 \frac{L_h^2}{B_h}.
\]
The uniform type estimate in Remark~\ref{rem:upper-type} now gives
\[
 \log\E Z_{n,d,h}(m_n)
 \le
 -\left(e^\omega-1+o(1)\right)
 \frac{nL_h^2}{B_h}
 +O_d(h\log n)
 \longrightarrow-\infty.
\]
Thus the pairing model has no dominating set of size exactly $m_n$ with high probability.  Any smaller dominating set can be padded to size $m_n$, so
\[
 \gamma_h(G)\ge m_n+1=\lceil T_n\rceil\ge T_n.
\]
Conditioning on simplicity transfers the conclusion to the uniformly random simple $d$-regular graph, and~\eqref{eq:two-path-dominates} gives the final assertion.
\end{proof}

\begin{proof}[Proof of Theorem~\ref{thm:random-near-critical}]
Set
\[
 T_n=\frac{nC_h^*}{B_h}.
\]
For large $n$, $C_h^*>0$.  If $T_n\le1$, the real-valued lower bound follows from $\gamma_h(G)\ge1$.  Otherwise put
\[
 m_n=\lceil T_n\rceil-1,
 \qquad
 C_n=\frac{m_nB_h}{n}.
\]
Then $m_n<T_n$ and $m_n/T_n\ge1/2$.  Choose $\eta>0$ such that $C_h^*\ge2\eta L_h$ eventually.  Then, for large $h$,
\[
 \eta L_h\le C_n\le C_h^*.
\]
Apply Theorem~\ref{thm:near-critical} with this $\eta$ and the same $W_h$.  Since $e^{-C_n}\ge e^{-C_h^*}=e^{W_h}L_h^2/B_h$ and $\frac12-o(1)\ge\frac13$ for sufficiently large $h$,
\[
 \log\E Z_{n,d,h}(m_n)
 \le-\frac13
 \frac{n e^{W_h}L_h^2}{B_h}
 +O_d(h\log n),
\]
which tends to $-\infty$ by~\eqref{eq:near-growth-condition}.  Markov's inequality shows that the configuration model has no dominating set of size exactly $m_n$ with high probability.  Any smaller dominating set could be padded to size $m_n$, so
\[
 \gamma_h(G)\ge m_n+1=\lceil T_n\rceil\ge T_n.
\]
Conditioning on simplicity transfers the conclusion to the uniformly random simple $d$-regular graph.
\end{proof}

\begin{proof}[Proof of Corollary~\ref{cor:random-fixed}]
Choose
\[
 W_h=\eps L_h-2\log L_h.
\]
Then $W_h\to\infty$,
\[
 C_h^*=(1-\eps)L_h,
 \qquad
 \frac{C_h^*}{L_h}=1-\eps>0,
\]
and
\[
 \frac{e^{W_h}L_h^2}{B_h}=B_h^{-(1-\eps)}.
\]
Thus Theorem~\ref{thm:random-near-critical} applies directly.
\end{proof}



\section{Numerical verification and conditioning}\label{sec:numerics}
The proofs above do not use numerical evidence.  The accompanying package nevertheless checks every stationary identity, the one-step shadowing expansions, the density targeting, and the activity law with controlled-precision arithmetic.

A cancellation issue discovered during numerical verification is important for reproducibility.  Direct evaluation of
\[
 \phi=\log Z_v-\frac d2\log Z_e
\]
can be ill-conditioned near criticality: both logarithms are large, and individual vertex terms of the form $S_i^d-(S_i-B_{i-1})^d$ may lose precision.  The consolidated solver evaluates these differences as
\[
 S_i^d\bigl[-\operatorname{expm1}(d\log(1-B_{i-1}/S_i))\bigr]
\]
and reports the microcanonical exponent through the corrected root-only formula~\eqref{eq:root-micro-correct}.  The direct partition-function route is retained as an independent comparison.  A row is rejected unless
\[
 \frac{|\Psi_{Z_v}-\Psi_{\rm root}|}{|\Psi_{\rm root}|}\le10^{-6}.
\]
Thus a cross-check residual comparable to the reported quantity gates the record rather than remaining hidden in an internal data structure.

For the diagnostic choice $W_h=\log\log B_h$, the regenerated cubic values are
\begin{center}
\begin{tabular}{cccc}
\toprule
$h$ & $C$ & $-\Psi/e^{-C}$ & activity-law ratio\\
\midrule
20 & $6.8451$ & $0.6484$ & $0.5873$\\
30 & $12.6345$ & $0.9418$ & $0.9520$\\
40 & $18.7408$ & $0.9757$ & $0.9956$\\
52 & $26.2980$ & $0.9819$ & $0.9997$\\
\bottomrule
\end{tabular}
\end{center}
The activity-law ratio is
\[
 \frac{-\log z-\log(1/\alpha)}{B_h(1-\alpha)^{B_h}}.
\]
The effective onset is slower when $C/\log B_h$ is small, consistent with uniformity only on compact intervals bounded away from zero.  Every numerical row records the checker hash, solver hash, precision, stationarity residual, telescoping residual, and the discrepancy between the direct and root-only free-energy routes.




The bounded-window audit evaluates both the true stationary optimizer and the capped-free profile at the center $s=0$.  The two quantities approach the coupon term from opposite sides:
\begin{center}
\begin{tabular}{cccc}
\toprule
$(d,h)$ & $(H-\Psi)/Q_h$ & cap loss$/Q_h$ & $C_h^{\rm ann}-\widehat C_{B_h}$\\
\midrule
$(3,20)$ & $0.8750056$ & $1.0593258$ & $-1.03077\times10^{-1}$\\
$(3,30)$ & $0.9902396$ & $1.0076650$ & $-8.00386\times10^{-3}$\\
$(3,40)$ & $0.9990005$ & $1.0009330$ & $-8.57353\times10^{-4}$\\
$(4,24)$ & $0.9999148$ & $1.0000702$ & $-6.78418\times10^{-5}$\\
$(4,32)$ & $0.9999988$ & $1.0000012$ & $-1.01763\times10^{-6}$\\
\bottomrule
\end{tabular}
\end{center}
The cap-loss and product identities are evaluated independently from the compact functional; their recorded residuals are below $10^{-120}$ in the displayed high-precision runs.  The stationary entropy-defect identity~\eqref{eq:defect-identity} and the graph-zero computations provide separate checks of the two sides of the proof.

\section{Consequences and open problems}\label{sec:open}
Theorems~\ref{thm:bounded-critical-window} and~\ref{thm:random-fixed-window} determine the annealed transition, its bounded window, and its universal scaling function.  The capped-free construction also isolates the mechanism behind the coupon term: at the critical scale, the leading entropy defect is the cost of forbidding one empty terminal branch set.  Two harder frontiers remain logically separate.

\paragraph{Quenched matching.}
The random-regular theorem is a first-moment lower bound.  A matching upper bound near $C_h^{\rm ann}$ would require proving that dominating sets actually exist above the annealed crossing.  Plausible routes include a second-moment analysis with an overlap variational problem, small-subgraph conditioning, or an algorithmic construction whose output reaches the scalar coupon coordinate.  The bounded-window theorem now supplies a precise target and separates any quenched correction from uncertainty in the annealed calculation.

\paragraph{Direct two-branch asymptotics.}
Internally two-path $(h,2)$ domination is stronger than ordinary distance domination, so the present lower bound transfers automatically.  It does not identify the parameter's own leading constant or critical correction.  A direct treatment must remember branch information without falsely treating two paths that later merge as internally disjoint.  Directed nonbacktracking cavity states are a natural exact intermediary, but the corresponding type system and variational reduction remain to be developed.

The exponent $1/7$ in the graph--scalar comparison is a bookkeeping exponent rather than a predicted optimum.  Improving it would sharpen finite-$h$ convergence but would not change the scaling function or any fixed-order inverse-logarithmic term.

\section*{Acknowledgements}
\emph{To be completed by the author before circulation.}

\begin{thebibliography}{99}
\bibitem{AignerFromme1984}
M. Aigner and M. Fromme,
\emph{A game of cops and robbers},
Discrete Applied Mathematics 8 (1984), 1--12; DOI: \href{https://doi.org/10.1016/0166-218X(84)90073-8}{10.1016/0166-218X(84)90073-8}.

\bibitem{CohenHonkalaLitsynLobstein1997}
G. Cohen, I. Honkala, S. Litsyn, and A. Lobstein,
\emph{Covering Codes},
North-Holland Mathematical Library 54, Elsevier, 1997.

\bibitem{CutlerRadcliffe2016}
J. Cutler and A. J. Radcliffe,
\emph{Counting dominating sets and related structures in graphs},
Discrete Mathematics 339 (2016), 1593--1599; DOI: \href{https://doi.org/10.1016/j.disc.2015.12.011}{10.1016/j.disc.2015.12.011}.

\bibitem{Duckworth2005}
W. Duckworth,
\emph{Randomized greedy algorithms for finding small $k$-dominating sets of regular graphs},
Random Structures \& Algorithms 27 (2005), 401--412; DOI: \href{https://doi.org/10.1002/rsa.20082}{10.1002/rsa.20082}.

\bibitem{DuckworthWormald2006}
W. Duckworth and N. C. Wormald,
\emph{On the independent domination number of random regular graphs},
Combinatorics, Probability and Computing 15 (2006), 513--522; DOI: \href{https://doi.org/10.1017/S0963548305007431}{10.1017/S0963548305007431}.

\bibitem{GlebovLiebenauSzabo2015}
R. Glebov, A. Liebenau, and T. Szab\'o,
\emph{On the concentration of the domination number of the random graph},
SIAM Journal on Discrete Mathematics 29 (2015), 1186--1206; DOI: \href{https://doi.org/10.1137/12090054X}{10.1137/12090054X}.

\bibitem{HabibullaQin2019}
Y. Habibulla and S.-m. Qin,
\emph{Two-distance minimal dominating set problem studied by statistical mechanics and simulated annealing},
arXiv:1910.07933, 2019; \href{https://arxiv.org/abs/1910.07933}{arXiv link}.

\bibitem{Janson2009}
S. Janson,
\emph{The probability that a random multigraph is simple},
Combinatorics, Probability and Computing 18 (2009), 205--225; DOI: \href{https://doi.org/10.1017/S0963548308009644}{10.1017/S0963548308009644}.

\bibitem{LuPeng2012}
L. Lu and X. Peng,
\emph{On Meyniel's conjecture of the cop number},
Journal of Graph Theory 71 (2012), 192--205; DOI: \href{https://doi.org/10.1002/jgt.20642}{10.1002/jgt.20642}.

\bibitem{Neuwirth2026Tube}
L. Neuwirth,
\emph{Branch-tube persistence and static coverage in tree-ball geometry},
preprint, 2026; \href{https://levineuwirth.org/essays/branch-based-local-capture-in-tree-balls/}{project page}.

\bibitem{PralatWormald2016}
P. Pra\l at and N. C. Wormald,
\emph{Meyniel's conjecture holds for random graphs},
Random Structures \& Algorithms 48 (2016), 396--421; DOI: \href{https://doi.org/10.1002/rsa.20587}{10.1002/rsa.20587}.

\bibitem{PralatWormald2019}
P. Pra\l at and N. C. Wormald,
\emph{Meyniel's conjecture holds for random $d$-regular graphs},
Random Structures \& Algorithms 55 (2019), 719--741; DOI: \href{https://doi.org/10.1002/rsa.20874}{10.1002/rsa.20874}.

\bibitem{RockafellarWets1998}
R. T. Rockafellar and R. J.-B. Wets,
\emph{Variational Analysis},
Grundlehren der mathematischen Wissenschaften 317, Springer, 1998; DOI: \href{https://doi.org/10.1007/978-3-642-02431-3}{10.1007/978-3-642-02431-3}.

\bibitem{ScottSudakov2011}
A. Scott and B. Sudakov,
\emph{A bound for the cops and robbers problem},
SIAM Journal on Discrete Mathematics 25 (2011), 1438--1442; DOI: \href{https://doi.org/10.1137/100812963}{10.1137/100812963}.

\bibitem{Wormald1999}
N. C. Wormald,
\emph{Models of random regular graphs},
in Surveys in Combinatorics, 1999, London Mathematical Society Lecture Note Series 267, Cambridge University Press, 1999, pp. 239--298; DOI: \href{https://doi.org/10.1017/CBO9780511721335.010}{10.1017/CBO9780511721335.010}.

\bibitem{ZhaoHabibullaZhou2015}
J.-H. Zhao, Y. Habibulla, and H.-J. Zhou,
\emph{Statistical mechanics of the minimum dominating set problem},
Journal of Statistical Physics 159 (2015), 1154--1174; DOI: \href{https://doi.org/10.1007/s10955-015-1220-2}{10.1007/s10955-015-1220-2}.
\end{thebibliography}

\end{document}
